%% GhostKitchen — Complete Engineering Journey, Architecture Evolution & Technical Project Book
%% March – June 2026
%% Compile: pdflatex (run twice for ToC/labels)

\documentclass[11pt,a4paper,twoside]{book}

\usepackage[T1]{fontenc}
\usepackage[utf8]{inputenc}
\usepackage[inner=2.8cm,outer=2.2cm,top=2.5cm,bottom=2.8cm]{geometry}
\usepackage{lmodern}
\usepackage{microtype}
\usepackage{booktabs}
\usepackage{longtable}
\usepackage{array}
\usepackage{colortbl}
\usepackage{xcolor}
\usepackage{hyperref}
\usepackage{enumitem}
\usepackage{titlesec}
\usepackage{fancyhdr}
\usepackage{tcolorbox}
\usepackage{tabularx}
\usepackage{multirow}
\usepackage{amsmath}
\usepackage{fontawesome5}
\usepackage{listings}
\usepackage{tikz}
\usepackage{pgfplots}
\usepackage{graphicx}
\usepackage{wrapfig}
\usepackage{mdframed}
\usepackage{caption}
\usepackage{subcaption}
\usepackage{lscape}
\usepackage{makecell}
\pgfplotsset{compat=1.18}
\usetikzlibrary{arrows.meta,positioning,shapes.geometric,fit,backgrounds,decorations.pathreplacing,calc}
\tcbuselibrary{skins,breakable,listings}

% ── Colour Palette ─────────────────────────────────────────────────────────────
\definecolor{brand}{HTML}{FF5200}
\definecolor{branddark}{HTML}{CC4100}
\definecolor{dark}{HTML}{111827}
\definecolor{muted}{HTML}{6B7280}
\definecolor{light}{HTML}{F9FAFB}
\definecolor{ok}{HTML}{059669}
\definecolor{warn}{HTML}{D97706}
\definecolor{crit}{HTML}{DC2626}
\definecolor{info}{HTML}{2563EB}
\definecolor{purple}{HTML}{7C3AED}
\definecolor{codebg}{HTML}{F3F4F6}
\definecolor{rowalt}{HTML}{FFF7F4}
\definecolor{headbg}{HTML}{FF5200}

\hypersetup{colorlinks=true,linkcolor=branddark,urlcolor=info,citecolor=brand,
  pdftitle={GhostKitchen Engineering Journey},pdfauthor={GhostKitchen Engineering}}

% ── Chapter / Section Style ────────────────────────────────────────────────────
\titleformat{\chapter}[display]
  {\normalfont\huge\bfseries\color{brand}}
  {\color{muted}\normalsize CHAPTER \thechapter}{12pt}{\Huge}[\vspace{4pt}\color{brand}\titlerule]
\titlespacing*{\chapter}{0pt}{-10pt}{20pt}
\titleformat{\section}{\large\bfseries\color{dark}}{\thesection.}{0.5em}{}[\vspace{-2pt}{\color{brand}\rule{\linewidth}{0.6pt}}]
\titleformat{\subsection}{\normalsize\bfseries\color{dark}}{\thesubsection}{0.5em}{}
\titleformat{\subsubsection}{\small\bfseries\itshape\color{dark}}{}{0em}{}

% ── Headers / Footers ──────────────────────────────────────────────────────────
\pagestyle{fancy}
\fancyhf{}
\renewcommand{\headrulewidth}{0.4pt}
\fancyhead[LE]{\small\color{muted}\leftmark}
\fancyhead[RO]{\small\color{muted}\rightmark}
\fancyhead[RE,LO]{\small\color{brand}\textbf{GhostKitchen}}
\fancyfoot[C]{\small\color{muted}\thepage}
\fancypagestyle{plain}{\fancyhf{}\fancyfoot[C]{\small\color{muted}\thepage}}

% ── Code Listings ──────────────────────────────────────────────────────────────
\lstset{
  backgroundcolor=\color{codebg},basicstyle=\ttfamily\footnotesize,
  breaklines=true,frame=single,framesep=4pt,rulecolor=\color{muted},
  keywordstyle=\color{info},commentstyle=\color{muted}\itshape,
  stringstyle=\color{ok},numbers=left,numberstyle=\tiny\color{muted},
  numbersep=6pt,xleftmargin=12pt,
}

% ── Coloured Table Helpers ─────────────────────────────────────────────────────
\newcommand{\th}[1]{\textbf{\color{white}#1}}
\newcommand{\cellok}[1]{\cellcolor{ok!15}\textbf{\color{ok}#1}}
\newcommand{\cellwarn}[1]{\cellcolor{warn!20}\textbf{\color{warn}#1}}
\newcommand{\cellcrit}[1]{\cellcolor{crit!15}\textbf{\color{crit}#1}}
\newcommand{\cellinfo}[1]{\cellcolor{info!10}\textbf{\color{info}#1}}
\newcommand{\cellmuted}[1]{\cellcolor{muted!12}\textit{\color{muted}#1}}
\newcommand{\yes}{\textcolor{ok}{\faCheckCircle}}
\newcommand{\no}{\textcolor{crit}{\faTimesCircle}}
\newcommand{\half}{\textcolor{warn}{\faAdjust}}

% ── tcolorbox Styles ──────────────────────────────────────────────────────────
\tcbset{
  insight/.style={colback=info!6,colframe=info,fonttitle=\bfseries,breakable},
  warn/.style={colback=warn!8,colframe=warn,fonttitle=\bfseries,breakable},
  crit/.style={colback=crit!6,colframe=crit,fonttitle=\bfseries,breakable},
  success/.style={colback=ok!6,colframe=ok,fonttitle=\bfseries,breakable},
  story/.style={colback=light,colframe=muted!40,fonttitle=\bfseries\color{brand},
    title=Engineering Decision,breakable},
  commit/.style={colback=brand!5,colframe=brand!50,fonttitle=\bfseries\footnotesize,breakable},
}

% ── Era Timeline Node ──────────────────────────────────────────────────────────
\newcommand{\erabox}[4]{%
  \begin{tcolorbox}[colback=#1!8,colframe=#1,title={\textbf{#2}},fonttitle=\bfseries\small]
  \textbf{Dates:} #3\\[2pt]#4
  \end{tcolorbox}
}

\begin{document}

% ╔══════════════════════════════════════════════════════╗
% ║                    FRONT MATTER                      ║
% ╚══════════════════════════════════════════════════════╝

\frontmatter

% ── Half-Title ────────────────────────────────────────────────────────────────
\begin{titlepage}
  \centering
  \vspace*{2cm}
  \begin{tikzpicture}
    \fill[brand] (0,0) rectangle (14,1.2);
    \node[text=white,font=\Huge\bfseries] at (7,0.6) {GhostKitchen};
  \end{tikzpicture}\\[0.6cm]
  {\LARGE\bfseries\color{dark} Complete Engineering Journey}\\[0.3cm]
  {\large\color{muted} Architecture Evolution \& Technical Project Book}\\[0.2cm]
  \begin{tikzpicture}
    \fill[brand] (0,0) rectangle (14,0.08);
  \end{tikzpicture}\\[2.5cm]

  \begin{tabular}{p{4cm}p{9cm}}
    \textbf{Platform}     & Multi-role food delivery — customers, restaurants, riders, admins \\[4pt]
    \textbf{Backend}      & Node.js 20 · Express · Prisma ORM · PostgreSQL · Redis · Socket.IO \\[4pt]
    \textbf{Frontend}     & Next.js 15 App Router · React · TailwindCSS · Zustand · Leaflet \\[4pt]
    \textbf{Payments}     & Cashfree PG (create-order → webhook → verify → refund) \\[4pt]
    \textbf{Hosting}      & Render (backend) · Vercel (frontend) · Neon/Supabase (PostgreSQL) \\[4pt]
    \textbf{Period}       & March 26, 2026 – June 14, 2026 (81 days, 194 commits) \\[4pt]
    \textbf{Scale}        & 49,955 source LOC · 11,613 test LOC · 1,036 automated tests \\[4pt]
    \textbf{Authors}      & GhostKitchen Engineering Team \\[4pt]
    \textbf{Version}      & 1.0 · June 2026 \\
  \end{tabular}

  \vfill
  \begin{tcolorbox}[colback=brand!6,colframe=brand,width=\textwidth]
  \centering
  \textbf{From prototype to production-ready platform in 81 days.}\\
  This book tells the complete engineering story — the decisions, the crises,
  the pivots, and the architecture that emerged.
  \end{tcolorbox}
  \vspace{0.5cm}
  {\small\color{muted}\textbf{CONFIDENTIAL — Internal Engineering Documentation}}
\end{titlepage}

\cleardoublepage
\tableofcontents
\cleardoublepage

% ── Preface ───────────────────────────────────────────────────────────────────
\chapter*{Preface}
\addcontentsline{toc}{chapter}{Preface}

GhostKitchen began as a prototype food-delivery application and evolved, over
eleven weeks, into a production-ready multi-role platform handling payments,
real-time GPS tracking, automated incident detection, and executive analytics.

This book documents that journey \emph{from source}: every architectural
decision is grounded in the commit history, every schema diagram derived from
the actual migration files, every benchmark drawn from the test suite.  It is
not a marketing document.  It is an engineering record.

The story has five acts:

\begin{enumerate}[leftmargin=1.5em]
  \item \textbf{The Prototype} (March 26 – April 7) — First working checkout flow,
        Cashfree payments, Socket.IO real-time tracking.
  \item \textbf{The Production Crisis} (April 7–10) — Twenty-three sequential bug-fix
        commits exposing every gap between a local demo and a live system.
  \item \textbf{The Great Rebuild} (June 3–7) — Complete architectural overhaul:
        multi-role RBAC, token-based auth replacing NextAuth, self-healing migrations.
  \item \textbf{Product Maturity} (June 10–13) — Admin console, 6,708-city onboarding,
        WCAG accessibility, security sprints, test coverage from 17.6\% to 76.6\%.
  \item \textbf{Operations Excellence} (June 14) — Live GPS map, incident engine,
        executive analytics, distributed locks, observability hardening.
\end{enumerate}

Each chapter that follows examines one dimension of this evolution in depth.
Architecture diagrams are drawn from the final system state; timelines reflect
the actual commit chronology.

\mainmatter

% ╔══════════════════════════════════════════════════════╗
% ║  CHAPTER 1 — PROJECT OVERVIEW                        ║
% ╚══════════════════════════════════════════════════════╝
\chapter{Project Overview}

\section{What Is GhostKitchen?}

GhostKitchen is a production-grade, multi-role food delivery platform built for
the Indian market.  It connects four distinct user types — customers, restaurant
owners, delivery riders, and platform administrators — through a single
application with role-gated dashboards, real-time order tracking, and integrated
payment processing via Cashfree PG.

The platform is \emph{not} a tutorial project.  It implements:

\begin{itemize}[leftmargin=1.4em]
  \item Hashed, rotating refresh tokens with concurrent-race protection
  \item Server-side socket room authorization backed by Prisma ownership queries
  \item HMAC-SHA256 webhook verification on raw request bodies
  \item Distributed Redis locks preventing duplicate cron-job execution across instances
  \item A three-tier observability stack: Winston structured logs → Sentry with release tags → enriched health endpoints
  \item An automated incident engine that detects, escalates, and notifies without human intervention
\end{itemize}

\section{Problem Statement}

Food delivery in India is dominated by aggregators (Swiggy, Zomato) that charge
15–30\% commission.  Independent restaurant operators need a platform they
control — one that can run on their own domain, integrate their existing payment
accounts, and provide analytics without a middleman.

GhostKitchen targets this operator-first model: the restaurant owner deploys
their own instance, retains all payment revenue, and gains operational
intelligence that aggregators keep proprietary.

\section{Business Goals}

\begin{enumerate}[leftmargin=1.5em]
  \item Enable restaurant owners to accept and manage orders without third-party
        commission.
  \item Provide customers with a smooth browse→cart→checkout→track experience
        comparable to national aggregators.
  \item Give delivery riders a transparent earnings dashboard with session tracking.
  \item Equip platform administrators with real-time operational intelligence,
        incident detection, and executive KPIs.
\end{enumerate}

\section{Target Users}

\begin{center}\small
\begin{tabular}{llp{6.5cm}}
\toprule
\rowcolor{headbg}\th{Role} & \th{Primary Device} & \th{Core Needs} \\
\midrule
Customer        & Mobile / web & Browse, search, order, track, re-order, reviews \\
\rowcolor{rowalt}Restaurant     & Tablet / desktop & Receive orders, manage menu, view analytics \\
Delivery Rider  & Mobile       & Accept assignments, GPS-stream position, track earnings \\
\rowcolor{rowalt}Admin          & Desktop          & Operations visibility, incident management, executive KPIs \\
\bottomrule
\end{tabular}
\end{center}

\section{Technical Goals}

\begin{tcolorbox}[insight,title=System Requirements]
\begin{itemize}[leftmargin=1.2em,nosep]
  \item \textbf{Sub-100ms API responses} for browse, search, and cart operations
  \item \textbf{Real-time delivery tracking} with GPS precision via Socket.IO
  \item \textbf{Zero-downtime deploys} via self-healing migration scripts
  \item \textbf{Multi-instance horizontal scaling} with shared Redis state
  \item \textbf{Full OWASP Top 10 coverage} across all API surfaces
  \item \textbf{WCAG 2.1 AA accessibility} for all customer-facing pages
\end{itemize}
\end{tcolorbox}

\section{Project Scale in Numbers}

\begin{center}\small
\begin{tabular}{lr lr}
\toprule
\textbf{Metric} & \textbf{Value} & \textbf{Metric} & \textbf{Value} \\
\midrule
Total commits        & 194          & Backend source LOC   & 22,767 \\
Development days     & 81           & Frontend source LOC  & 27,188 \\
Database models      & 20           & Test LOC             & 11,613 \\
Prisma migrations    & 21           & Total tests          & 1,036  \\
API endpoints        & 89           & Test files           & 42     \\
Frontend pages       & 44           & npm dependencies     & 49     \\
Socket event types   & 14           & Env variables        & 18     \\
\bottomrule
\end{tabular}
\end{center}

% ╔══════════════════════════════════════════════════════╗
% ║  CHAPTER 2 — THE FIVE ENGINEERING ERAS               ║
% ╚══════════════════════════════════════════════════════╝
\chapter{The Five Engineering Eras}

\section{Overview}

GhostKitchen's 194 commits fall into five distinct engineering eras, each with
a dominant theme, a defining challenge, and measurable outcomes.

\begin{figure}[h!]\centering
\begin{tikzpicture}[
  era/.style={rectangle,rounded corners=4pt,minimum height=1.1cm,
    text width=2.4cm,align=center,font=\footnotesize\bfseries},
  line/.style={-Stealth,thick,color=muted}
]
\node[era,fill=info!20,draw=info]       (e1) at (0,0)   {Era I\\Genesis\\Mar 26–Apr 7};
\node[era,fill=crit!15,draw=crit]       (e2) at (3,0)   {Era II\\Crisis\\Apr 7–10};
\node[era,fill=purple!15,draw=purple]   (e3) at (6,0)   {Era III\\Rebuild\\Jun 3–7};
\node[era,fill=warn!20,draw=warn]       (e4) at (9,0)   {Era IV\\Maturity\\Jun 10–13};
\node[era,fill=ok!15,draw=ok]           (e5) at (12,0)  {Era V\\Excellence\\Jun 14};
\draw[line] (e1) -- (e2);
\draw[line] (e2) -- (e3);
\draw[line] (e3) -- (e4);
\draw[line] (e4) -- (e5);
\node[font=\tiny,color=muted,below=0.3cm of e1] {47 commits};
\node[font=\tiny,color=muted,below=0.3cm of e2] {23 commits};
\node[font=\tiny,color=muted,below=0.3cm of e3] {52 commits};
\node[font=\tiny,color=muted,below=0.3cm of e4] {42 commits};
\node[font=\tiny,color=muted,below=0.3cm of e5] {30 commits};
\end{tikzpicture}
\caption{GhostKitchen's five engineering eras}
\end{figure}

\section{Era I — Genesis (March 26 – April 7, 47 commits)}

\subsection{What Was Built}

The first commit on March 26 established a basic Express backend with restaurant
browsing, a cart system, and Cashfree payment integration.  The frontend was a
Next.js 14 application using NextAuth for session management.

\begin{tcolorbox}[commit,title=Milestone commits]
\begin{itemize}[leftmargin=1em,nosep,topsep=2pt]
  \item \texttt{cad7cd2} — Phase 1: Production-Ready Auth System (Apr 9)
  \item \texttt{8aaa7ff} — Phase 2: Cart system with DB persistence (Apr 9)
  \item \texttt{b36822c} — Phase 3: Order System — Cart→Checkout→Order (Apr 9)
  \item \texttt{b6a92f1} — Cashfree Payment Integration (Apr 9)
  \item \texttt{b8b7fe4} — Phase 2: Real-time tracking with Socket.IO (Apr 10)
  \item \texttt{c13d76b} — Phase 3: Admin panel and dashboard (Apr 10)
\end{itemize}
\end{tcolorbox}

\subsection{Architecture at Day 1}

The initial stack was deliberately minimal: a single Express server, Prisma
connected to a Neon PostgreSQL database, no Redis, and Socket.IO in in-memory
mode.  Authentication used NextAuth on the frontend passing session cookies to
the backend — a pattern that would prove catastrophic in the cross-origin
Vercel/Render deployment.

\subsection{Initial Limitations}

\begin{itemize}[leftmargin=1.4em]
  \item \textbf{Single-role model}: users were either customers or admins; no restaurant owner or rider roles
  \item \textbf{No refresh token rotation}: JWT access tokens were long-lived
  \item \textbf{No webhook idempotency}: duplicate Cashfree webhooks created duplicate orders
  \item \textbf{No rate limiting}: all endpoints were open to brute force
  \item \textbf{CORS wildcards}: \texttt{Access-Control-Allow-Origin: *} in early config
\end{itemize}

\section{Era II — The Production Crisis (April 7–10, 23 commits)}

\subsection{The Problem}

When the prototype was deployed to Render + Vercel, it produced a cascade of
failures.  The commit log for April 7–10 reads as a triage log:

\begin{lstlisting}[language=bash,numbers=none]
ISSUE 5: Add transaction locking for race condition prevention
ISSUE 6: Fix order status override with conditional update
ISSUE 7: Optimize cron job with single atomic updateMany
ISSUE 8: Add database indexes for Order model queries
ISSUE 9: Add structured logging context to payment service
FINAL PRODUCTION GAPS: Fix 3 critical vulnerabilities
\end{lstlisting}

\subsection{Root Causes}

The crisis revealed that local development masked critical production gaps:

\begin{center}\small
\begin{tabular}{p{5cm}p{7.5cm}}
\toprule
\rowcolor{headbg}\th{Symptom} & \th{Root Cause} \\
\midrule
Restaurant API 500 & Missing columns in production DB (slug, lat, lng) \\
\rowcolor{rowalt}Duplicate orders & Cashfree webhook had no idempotency guard \\
Checkout double-submit & No server-side duplicate check on order creation \\
\rowcolor{rowalt}CORS failures & Hardcoded origin list missed Vercel preview URLs \\
Payment race condition & Concurrent webhooks both passed SELECT-then-UPDATE \\
\rowcolor{rowalt}Slow order queries & No indexes on \texttt{Order.restaurantId}, \texttt{Order.customerId} \\
\bottomrule
\end{tabular}
\end{center}

\subsection{What the Crisis Taught}

\begin{tcolorbox}[story]
This 72-hour debugging marathon produced a discipline that would govern every
subsequent phase: never ship a feature without thinking through the production
deployment, the database state, and the concurrent-request scenario.  The
23 rapid-fire commits of Era II became the platform's first real stress test.
\end{tcolorbox}

\section{Era III — The Great Rebuild (June 3–7, 52 commits)}

\subsection{The Architectural Decision}

Between April 10 and June 3 — a 54-day gap with no commits — the team made a
fundamental decision: the NextAuth-based single-role architecture could not be
extended to support restaurant owners and delivery riders without a rewrite.
The rebuild introduced:

\begin{itemize}[leftmargin=1.4em]
  \item \textbf{Multi-role JWT tokens}: a single user can hold CUSTOMER, RESTAURANT, DELIVERY, and ADMIN roles simultaneously
  \item \textbf{Role switching}: token re-issue endpoint with active-role encoding
  \item \textbf{Authorization header auth}: replaced cookie-based auth to circumvent Safari third-party cookie blocking (Vercel ↔ Render cross-origin)
  \item \textbf{Zustand state management}: replaced NextAuth sessions with a persisted client-side store
  \item \textbf{Self-healing migrations}: idempotent DDL (\texttt{ADD COLUMN IF NOT EXISTS}) to patch production without downtime
\end{itemize}

\subsection{The Cookie Crisis (June 4–5)}

The rebuild triggered a new crisis: the frontend's Zustand auth store would
hydrate from localStorage, call \texttt{/auth/me} with an Authorization header,
then immediately clear state on a 401 — causing an infinite re-render loop.
Six commits on June 4 tackled this in sequence:

\begin{tcolorbox}[commit,title=Auth loop resolution (June 4--5)]
\begin{itemize}[leftmargin=1em,nosep,topsep=2pt]
  \item \texttt{bef7fcf} — Wait for Zustand hydration before auth redirect
  \item \texttt{15d1a98} — Don't clear auth state on cross-origin 401
  \item \texttt{24f47a6} — Use Authorization header to bypass third-party cookie blocking
  \item \texttt{17a035e} — Allow Authorization header in preflight CORS response
\end{itemize}
\end{tcolorbox}

\subsection{Restaurant Healing Logic}

A subtle production issue emerged when restaurant owners registered their
restaurant and then lost their \texttt{restaurantId} claim on token refresh.
The fix introduced an \emph{auto-heal} mechanism on \texttt{/auth/me}: if the
user has a RESTAURANT role but no \texttt{restaurantId} in their token, the
server queries the DB and re-issues a corrected token.

\section{Era IV — Product Maturity (June 10–13, 42 commits)}

\subsection{Feature Explosion}

Once the architecture stabilized, June 10–13 delivered features at pace:

\begin{center}\footnotesize
\begin{longtable}{p{3.5cm}p{9cm}}
\toprule
\rowcolor{headbg}\th{Date} & \th{Commits / Features} \\
\midrule
June 10 & Admin DMS drawer, \texttt{/bootstrap-admin}, city-search onboarding, 6,708 Indian cities, dark admin shell redesign, settings enforcement, user management (suspend/block/delete/role-change) \\
\rowcolor{rowalt}June 11 & Favorites UI, smart homepage, address-based checkout, customer refunds, field-name audit sweep, backend hardening sprint, Sentry integration, CI/CD pipeline \\
June 12 & ETA engine, recommendations engine, admin completions, launch-readiness security fixes (OOM guard, refund race, requireApproval), IDOR closure, WCAG 2.1 AA fixes, 103-test suite \\
\rowcolor{rowalt}June 13 & Backend coverage 17.6\%→76.6\%, CI/CD repair, README overhaul, maintenance mode, self-healing deploy P3009 fix \\
\bottomrule
\end{longtable}
\end{center}

\subsection{The 6{,}708-City Dataset}

\begin{tcolorbox}[story]
Restaurant onboarding required a city-selection field.  The initial implementation
used a handful of hardcoded major cities, which immediately broke for any
restaurant outside the top 10 metros.  The solution was to embed an open-source
dataset of all 6,708 Indian cities and towns into the frontend bundle, paired with
a fuzzy-search autocomplete.  Commit \texttt{4216492} added this dataset; the
bundle impact was acceptable because Next.js code-splits it to the onboarding
page only.
\end{tcolorbox}

\section{Era V — Operations Excellence (June 14, 30 commits)}

June 14 was the most productive single day in the project: 30 commits delivering
five major systems.

\begin{center}\small
\begin{tabular}{p{3.5cm}p{9cm}}
\toprule
\rowcolor{headbg}\th{System} & \th{Description} \\
\midrule
Live Operations Map      & Real-time rider GPS on Leaflet with heatmap, clustering, route display \\
\rowcolor{rowalt}Ops Intelligence Waves A–C & Incident engine, scoring, leaderboards, executive analytics \\
Security Sprints 1–2     & OWASP hardening, Permissions-Policy, audit trails, regression suites \\
\rowcolor{rowalt}Distributed Locks         & Redis SETNX guards on all three cron jobs \\
Observability             & JSON logs, Sentry release tags, X-Request-ID, enriched health endpoints \\
\bottomrule
\end{tabular}
\end{center}

% ╔══════════════════════════════════════════════════════╗
% ║  CHAPTER 3 — DATABASE EVOLUTION                      ║
% ╚══════════════════════════════════════════════════════╝
\chapter{Database Evolution}

\section{Migration Philosophy}

GhostKitchen's database evolved through 21 migrations over 11 weeks.  From
migration 5 onwards, all DDL used \texttt{IF NOT EXISTS} guards, making every
migration safe to replay — critical for the self-healing deployment pipeline
that had to recover from P3009 migration-lock failures on Render.

\section{Migration Chronology}

\begin{center}\footnotesize
\begin{longtable}{p{5.5cm}p{7.5cm}}
\toprule
\rowcolor{headbg}\th{Migration} & \th{Purpose \& Key Changes} \\
\midrule
\texttt{20260603\_baseline\_multi\_role} &
  Foundation schema. Defines 4 enums (Role, OrderStatus, PaymentStatus, DiscountType).
  Creates User, RefreshToken, Restaurant, MenuItem, Order, CartItem, Review,
  Coupon, Payment, PaymentWebhook. \\
\rowcolor{rowalt}\texttt{20260603\_refresh\_token\_hash} &
  Security: renames \texttt{RefreshToken.token} → \texttt{tokenHash} to clarify
  the stored value is a SHA-256 hash, not a raw secret. \\
\texttt{20260603\_security\_hardening} &
  Adds AuditLog table with indexes on userId, action, createdAt.
  Adds composite index on MenuItem(restaurantId, isAvailable) for menu queries. \\
\rowcolor{rowalt}\texttt{20260604\_add\_missing\_columns} &
  Production repair: adds slug, lat, lng, isOpen, deliveryRadius, rating,
  reviewCount, cuisine, minOrder, deliveryFee, estimatedDelivery to Restaurant
  using IF NOT EXISTS to avoid failures on partially-upgraded instances. \\
\texttt{20260604\_add\_user\_roles\_column} &
  Multi-role: adds \texttt{roles Role[] DEFAULT ARRAY['CUSTOMER']} to User,
  enabling one account to hold multiple platform roles simultaneously. \\
\rowcolor{rowalt}\texttt{20260604\_rename\_refresh\_token\_column} &
  Idempotent recreation of RefreshToken table with correct schema + composite
  index on (userId, expiresAt) for token cleanup queries. \\
\texttt{20260605\_complete\_systems} &
  Adds reviewCount, suspended to Restaurant. Adds isActive, description,
  restaurantId to Coupon. Creates Notification table for in-app alerts. \\
\rowcolor{rowalt}\texttt{20260607\_fix\_menuitem\_columns} &
  Adds isVeg, isBestSeller, sortOrder, isAvailable to MenuItem.
  Adds category, description indexes for menu filtering. \\
\texttt{20260610\_add\_user\_is\_blocked} &
  Security: adds isBlocked flag enabling admin to lock out accounts. \\
\rowcolor{rowalt}\texttt{20260610\_add\_restaurant\_role\_value} &
  Enum migration: adds RESTAURANT value to the Role enum (was missing). \\
\texttt{20260610\_add\_user\_is\_suspended} &
  Operations: adds isSuspended to distinguish temporary suspension from hard block. \\
\rowcolor{rowalt}\texttt{20260610\_site\_config\_restaurant\_approved} &
  Platform governance: creates singleton SiteConfig table (maintenance mode,
  feature flags, fee settings). Adds isApproved to Restaurant. \\
\texttt{20260610\_status\_note\_maintenance\_reason} &
  UX: adds statusNote to Restaurant for "back at 5pm" messaging.
  Adds maintenanceReason to SiteConfig. \\
\rowcolor{rowalt}\texttt{20260610\_maintenance\_scheduling} &
  Operations: adds maintenanceScheduledAt and maintenanceEndsAt to SiteConfig,
  enabling cron-driven maintenance window automation. \\
\texttt{20260610\_order\_analytics\_index} &
  Performance: composite index on Order(restaurantId, createdAt) accelerating
  restaurant analytics queries by 10–50× on large datasets. \\
\rowcolor{rowalt}\texttt{20260611\_addresses\_refunds\_favorites} &
  Customer features: creates Address (multi-address with default flag),
  Favorite (user↔restaurant unique constraint), Refund tables. \\
\texttt{20260611\_add\_order\_estimated\_delivery} &
  ETA: adds estimatedDelivery timestamp to Order, populated by the ETA engine
  at order confirmation. \\
\rowcolor{rowalt}\texttt{20260613\_add\_rider\_location} &
  GPS: creates RiderLocation with unique index on riderId and spatial index on
  (latitude, longitude) for proximity queries. \\
\texttt{20260614\_add\_incident} &
  Ops Intelligence: creates Incident with IncidentSeverity/Status enums.
  Indexes on status, severity, createdAt for dashboard queries. \\
\rowcolor{rowalt}\texttt{20260614\_incident\_escalation} &
  Wave B: adds acknowledgedAt, escalatedAt, escalationCount, lastNotifiedAt
  to Incident. Creates IncidentEvent for immutable audit timeline. \\
\texttt{20260614\_order\_lifecycle\_rider\_sessions} &
  Delivery tracking: adds acceptedAt, readyAt, pickedUpAt timestamps to Order.
  Creates RiderSession for online-time and earnings analytics. \\
\bottomrule
\end{longtable}
\end{center}

\section{Database Model Map}

\begin{figure}[h!]\centering
\begin{tikzpicture}[
  model/.style={rectangle,rounded corners=3pt,draw=dark,fill=light,
    minimum width=2.2cm,minimum height=0.65cm,font=\footnotesize\bfseries,align=center},
  rel/.style={-Stealth,color=muted,thin},
  group/.style={rectangle,rounded corners=8pt,draw=muted!40,dashed,
    fill=muted!5,inner sep=8pt},
]
% Auth cluster
\node[model,fill=info!15,draw=info] (user) at (0,8) {User};
\node[model,fill=info!8] (rt) at (3,8) {RefreshToken};
\node[model,fill=info!8] (session) at (3,6.5) {RiderSession};
\node[model,fill=info!8] (audit) at (0,6.5) {AuditLog};

% Core order cluster
\node[model,fill=brand!15,draw=brand] (order) at (7,7) {Order};
\node[model,fill=brand!8] (cart) at (5,5.5) {CartItem};
\node[model,fill=brand!8] (payment) at (9,5.5) {Payment};
\node[model,fill=brand!8] (webhook) at (11,7) {PaymentWebhook};
\node[model,fill=brand!8] (refund) at (9,8.5) {Refund};

% Content cluster
\node[model,fill=ok!15,draw=ok] (rest) at (0,3) {Restaurant};
\node[model,fill=ok!8] (menu) at (3,3) {MenuItem};
\node[model,fill=ok!8] (review) at (0,1.5) {Review};
\node[model,fill=ok!8] (coupon) at (3,1.5) {Coupon};
\node[model,fill=ok!8] (fav) at (6,3) {Favorite};
\node[model,fill=ok!8] (addr) at (6,1.5) {Address};
\node[model,fill=ok!8] (notif) at (5,8.5) {Notification};

% Ops cluster
\node[model,fill=purple!15,draw=purple] (incident) at (13,3.5) {Incident};
\node[model,fill=purple!8] (inevent) at (13,2) {IncidentEvent};
\node[model,fill=purple!8] (riderloc) at (11,3.5) {RiderLocation};

% Config
\node[model,fill=warn!15,draw=warn] (config) at (13,7) {SiteConfig};

% Relations
\draw[rel] (user) -- (rt);
\draw[rel] (user) -- (session);
\draw[rel] (user) -- (audit);
\draw[rel] (user) -- (order);
\draw[rel] (user) -- (cart);
\draw[rel] (user) -- (fav);
\draw[rel] (user) -- (addr);
\draw[rel] (user) -- (notif);
\draw[rel] (rest) -- (menu);
\draw[rel] (rest) -- (review);
\draw[rel] (rest) -- (order);
\draw[rel] (order) -- (payment);
\draw[rel] (order) -- (refund);
\draw[rel] (payment) -- (webhook);
\draw[rel] (incident) -- (inevent);
\draw[rel] (user) -- (riderloc);
\draw[rel] (order) -- (cart);
\end{tikzpicture}
\caption{Simplified GhostKitchen data model (20 models, 21 migrations)}
\end{figure}

\section{Schema Scale}

\begin{center}\small
\begin{tabular}{lrl lrl}
\toprule
\textbf{Model} & \textbf{Fields} && \textbf{Model} & \textbf{Fields} \\
\midrule
User          & 31 && Order     & 24 \\
Restaurant    & 23 && Incident  & 20 \\
SiteConfig    & 17 && MenuItem  & 15 \\
Payment       & 12 && Address   & 14 \\
Refund        & 9  && Notification & 9 \\
AuditLog      & 9  && RiderLocation & 8 \\
\bottomrule
\end{tabular}
\end{center}

The User model's 31 fields reflect the platform's multi-role complexity:
identity (name, email, passwordHash), geolocation (currentLat, currentLng),
availability (isAvailable for riders), moderation (isBlocked, isSuspended),
and role state (roles[], activeRole, restaurantId).

% ╔══════════════════════════════════════════════════════╗
% ║  CHAPTER 4 — BACKEND ARCHITECTURE                    ║
% ╚══════════════════════════════════════════════════════╝
\chapter{Backend Architecture}

\section{Module Structure}

The backend is a Node.js 20 / Express application structured around 15 feature
modules, each encapsulating routes, controller, and service.  The only shared
infrastructure is in \texttt{src/config/}, \texttt{src/middlewares/},
\texttt{src/utils/}, and \texttt{src/socket/}.

\begin{figure}[h!]\centering
\begin{tikzpicture}[
  box/.style={rectangle,draw=dark,fill=light,rounded corners=3pt,
    minimum width=2.8cm,minimum height=0.6cm,font=\footnotesize,align=center},
  infra/.style={rectangle,draw=info,fill=info!10,rounded corners=3pt,
    minimum width=2.8cm,minimum height=0.6cm,font=\footnotesize\bfseries,align=center},
]
% Entry
\node[infra] (srv) at (6,10) {server.js};
\node[infra] (app) at (6,8.8) {app.js (Express)};

% Middleware stack
\node[box] (m1) at (1.5,7.5) {Helmet/CSP};
\node[box] (m2) at (4.2,7.5) {CORS};
\node[box] (m3) at (6.8,7.5) {Rate Limiter};
\node[box] (m4) at (9.5,7.5) {Request Tracing};
\node[box] (m5) at (12,7.5) {Sanitize};

% Modules
\node[box,fill=brand!10,draw=brand] (auth) at (0,5.8) {auth};
\node[box] (cart) at (2.5,5.8) {cart};
\node[box] (orders) at (5,5.8) {orders};
\node[box] (payment) at (7.5,5.8) {payment};
\node[box] (rest) at (10,5.8) {restaurant};
\node[box] (del) at (12.5,5.8) {delivery};
\node[box] (admin) at (0,4.5) {admin};
\node[box] (ops) at (2.5,4.5) {ops};
\node[box] (exec) at (5,4.5) {exec};
\node[box] (notif) at (7.5,4.5) {notification};
\node[box] (user) at (10,4.5) {user};
\node[box] (coupon) at (12.5,4.5) {coupon};

% Infra
\node[infra] (prisma) at (2,2.8) {Prisma ORM};
\node[infra] (redis) at (5.5,2.8) {Redis (ioredis/Upstash)};
\node[infra] (socket) at (9,2.8) {Socket.IO};
\node[infra] (sentry) at (12,2.8) {Sentry};

% DB
\node[infra,fill=ok!15,draw=ok] (pg) at (6,1.5) {PostgreSQL (Neon)};

\draw[-Stealth,thick,color=brand] (srv) -- (app);
\draw[color=muted,thin] (app) -- (m1); \draw[color=muted,thin] (app) -- (m4);
\draw[color=muted,thin] (app) -- (m3); \draw[color=muted,thin] (app) -- (m5);
\draw[color=muted,thin] (auth) -- (prisma); \draw[color=muted,thin] (orders) -- (prisma);
\draw[color=muted,thin] (prisma) -- (pg);
\end{tikzpicture}
\caption{Backend layered architecture}
\end{figure}

\section{Request Lifecycle}

Every HTTP request flows through a deterministic pipeline:

\begin{enumerate}[leftmargin=1.5em]
  \item \textbf{Helmet} sets 10+ security headers (CSP, HSTS, X-Frame-Options).
  \item \textbf{Permissions-Policy} restricts camera, microphone, geolocation, payment.
  \item \textbf{Request Tracing} assigns \texttt{req.id = randomUUID()} and sets \texttt{X-Request-ID} response header.
  \item \textbf{CORS} validates origin against allowlist + Vercel preview pattern.
  \item \textbf{Compression} applies gzip at level 6 for API responses.
  \item \textbf{Body parsing} limits payloads to 50 KB.
  \item \textbf{XSS sanitisation} deep-cleans request bodies with the \texttt{xss} library.
  \item \textbf{Rate limiters} (tiered: general 200/15min, auth 10/15min, payment 20/min).
  \item \textbf{Route handler} calls authenticate → authorize → controller → service.
  \item \textbf{Error handler} classifies operational vs non-operational errors; routes to Sentry.
\end{enumerate}

\section{API Surface}

\begin{center}\small
\begin{longtable}{p{2.2cm}p{2cm}p{8.5cm}}
\toprule
\rowcolor{headbg}\th{Module} & \th{Endpoints} & \th{Key Operations} \\
\midrule
auth        & 6  & register, login, me, refresh, logout, logout-all \\
\rowcolor{rowalt}restaurant  & 14 & list, get, create, update, menu CRUD, analytics, trending, recommendations \\
orders      & 5  & list, get, place, update-status, reorder \\
\rowcolor{rowalt}payment     & 11 & create-order, verify, webhook, history, retry, refunds (customer+admin) \\
cart        & 5  & get, add, update, remove, clear \\
\rowcolor{rowalt}delivery    & 4  & accept-order, location ping, status toggle, earnings \\
admin       & 37 & full platform management across users, restaurants, orders, coupons, analytics \\
\rowcolor{rowalt}ops         & 11 & incidents, SLA, performance, leaderboards, alerts \\
exec        & 6  & KPIs, fleet, customers, restaurants, trends, summary \\
\rowcolor{rowalt}notification & 5 & list, read, read-all, delete, delete-all \\
user        & 9  & addresses, favorites, profile \\
\rowcolor{rowalt}role        & 7  & switch, register-restaurant, register-rider, add-restaurant, my-restaurant \\
review      & 6  & create, update, delete, get by order/restaurant, can-review \\
\rowcolor{rowalt}coupon      & 7  & validate, active, available, admin CRUD \\
config      & —  & getSiteConfigCached, updateSiteConfig (internal service) \\
\midrule
\textbf{Total} & \textbf{133} & \textbf{89 public endpoints + 44 internal routes} \\
\bottomrule
\end{longtable}
\end{center}

\section{Authentication System}

\subsection{Token Architecture}

GhostKitchen uses a dual-token model.  Access tokens are short-lived JWTs (15
minutes by default) encoding the user's roles, activeRole, and restaurantId.
Refresh tokens are long-lived (7 days) but stored only as SHA-256 hashes in
the database — the raw token is transmitted once and never stored in cleartext.

\begin{lstlisting}[language=JavaScript,caption={JWT payload structure}]
{
  userId: "cuid2_user_id",
  roles: ["CUSTOMER", "RESTAURANT"],
  activeRole: "RESTAURANT",
  restaurantId: "cuid2_restaurant_id",
  secondRole: "CUSTOMER",
  iat: 1718323200,
  exp: 1718324100  // 15 minutes
}
\end{lstlisting}

\subsection{Refresh Token Rotation}

The refresh endpoint implements rotation with a concurrent-race guard:

\begin{lstlisting}[language=JavaScript,caption={Rotation with race guard (auth.service.js)}]
// Rotate: delete old token atomically
const deleted = await prisma.refreshToken.deleteMany({ where: { tokenHash } });
if (deleted.count === 0) {
  // Another concurrent request already rotated this token — treat as revoked
  throw new AppError("Refresh token revoked or expired", 401);
}
// Issue new pair
const { accessToken, refreshToken } = generateTokenPair(buildTokenPayload(user));
await storeRefreshToken(user.id, refreshToken);
\end{lstlisting}

If two requests arrive simultaneously with the same refresh token, one will find
\texttt{deleted.count === 0} and be rejected, preventing token replay attacks.

\section{Authorization Architecture}

\subsection{Three-Layer RBAC}

\begin{center}\small
\begin{tabular}{p{3cm}p{9.5cm}}
\toprule
\rowcolor{headbg}\th{Layer} & \th{Mechanism} \\
\midrule
Router-level  & \texttt{router.use(authenticate, authorize("ADMIN"))} — blocks entire module \\
\rowcolor{rowalt}Route-level   & \texttt{roleMiddleware(["CUSTOMER"])} — per-endpoint role gate \\
Service-level & \texttt{getRestaurantByIdAndOwner(id, userId)} — ownership verification \\
\bottomrule
\end{tabular}
\end{center}

\subsection{IDOR Prevention}

Every resource with a per-user ownership invariant enforces it at the service
layer, independent of the role check.  Representative patterns:

\begin{lstlisting}[language=JavaScript,caption={Order IDOR gate (orders.controller.js)}]
async function canAccessOrder(order, user) {
  if (user.role === "ADMIN") return true;
  if (order.customerId === user.userId) return true;
  if (order.agentId && order.agentId === user.userId) return true;
  if (user.roles?.includes("RESTAURANT")) {
    const owned = await prisma.restaurant.findFirst({
      where: { id: order.restaurantId, ownerId: user.userId },
    });
    return !!owned;
  }
  return false;
}
\end{lstlisting}

% ╔══════════════════════════════════════════════════════╗
% ║  CHAPTER 5 — CUSTOMER PLATFORM                       ║
% ╚══════════════════════════════════════════════════════╝
\chapter{Customer Platform}

\section{Discovery \& Search}

\subsection{Restaurant Listing}

The primary browse experience presents restaurants with filtering by cuisine,
city, minimum order, delivery fee, and availability.  The \texttt{listRestaurants}
service supports cursor-based pagination and constructs a dynamic Prisma
\texttt{where} clause from query parameters.

\subsection{Search}

Search operates on restaurant name, cuisine, and description fields using
Prisma's case-insensitive \texttt{contains} mode.  A full-text index was
considered but deferred; at current scale the ILIKE query returns in under 20ms
with the existing btree index on \texttt{Restaurant.name}.

\subsection{Recommendations Engine}

Commit \texttt{59ff07a} introduced a personalised recommendations endpoint.
The algorithm ranks restaurants by a weighted composite of the user's own order
history and global popularity (review count × rating).  For authenticated users,
restaurants they've ordered from before receive a boost; for guests, the ranking
falls back to pure popularity.

\subsection{Trending}

A lightweight trending query aggregates order counts over the last 7 days,
returning the top 10 restaurants ordered by recent order volume.

\section{Cart System}

The cart is database-persisted, not session-based.  This decision was made
early (commit \texttt{8aaa7ff}) to support cross-device continuity.  CartItems
are keyed to the authenticated user and include a unique constraint on
(userId, menuItemId) to prevent duplicate additions.

\begin{tcolorbox}[warn,title=Cart Cross-Restaurant Guard]
The cart service enforces a single-restaurant constraint: adding an item from
a different restaurant clears the existing cart automatically.  This prevents
an order from containing items with incompatible delivery logistics.
\end{tcolorbox}

\section{Checkout \& Order Placement}

Order placement is a transactional operation that:

\begin{enumerate}[leftmargin=1.5em]
  \item Validates all CartItems still exist and are available
  \item Re-fetches prices from the database (never trusts client totals)
  \item Computes subtotal, delivery fee, and coupon discount server-side
  \item Verifies coupon eligibility (active, not expired, usage limit, minimum order)
  \item Atomically decrements \texttt{coupon.usedCount} and creates the Order in a Prisma transaction
  \item Emits \texttt{order:new} to the restaurant's Socket.IO room
  \item Sends an order-confirmation email via Resend
\end{enumerate}

\section{Payments}

\subsection{The Cashfree Integration}

GhostKitchen uses Cashfree PG's Node.js SDK for the payment lifecycle:

\begin{center}\small
\begin{tabular}{llp{7cm}}
\toprule
\rowcolor{headbg}\th{Step} & \th{Endpoint} & \th{Action} \\
\midrule
1. Create order  & \texttt{POST /payments/create-order} & \texttt{PGCreateOrder} — creates Cashfree order, returns \texttt{payment\_session\_id} \\
\rowcolor{rowalt}2. Checkout      & Frontend          & \texttt{cashfree.checkout()} loads Cashfree drop-in UI \\
3. Callback      & \texttt{/checkout/callback} & Reads \texttt{?order\_id} from URL, displays pending state \\
\rowcolor{rowalt}4. Webhook       & \texttt{POST /payments/webhook} & HMAC-verified, idempotent order confirmation \\
5. Verify        & \texttt{GET /payments/verify/:orderId} & Client-side poll after callback \\
\bottomrule
\end{tabular}
\end{center}

\subsection{Webhook Idempotency}

A critical bug in Era II was duplicate order creation from duplicate webhooks.
The fix uses \texttt{PaymentWebhook} as an idempotency table:

\begin{lstlisting}[language=JavaScript,caption={Webhook idempotency guard}]
const existing = await prisma.paymentWebhook.findUnique({
  where: { eventId: `${orderId}_${paymentStatus}_${cfPaymentId}` },
});
if (existing) {
  logger.info("Duplicate webhook — skipping", { orderId });
  return { duplicate: true };
}
\end{lstlisting}

\section{Order Tracking}

When an order is placed, the customer joins the Socket.IO room \texttt{order-\{id\}}.
Status changes from the restaurant (\texttt{CONFIRMED}, \texttt{PREPARING},
\texttt{OUT\_FOR\_DELIVERY}, \texttt{DELIVERED}) are emitted to this room via
\texttt{emitOrderStatusUpdated}.  When the rider goes out for delivery, their
GPS coordinates are relayed to the same room as \texttt{agent:location} events.

\section{Favorites, Addresses \& Reviews}

\begin{itemize}[leftmargin=1.4em]
  \item \textbf{Favorites}: toggle endpoints with a unique constraint on (userId, restaurantId); the homepage surfaces favourited restaurants first.
  \item \textbf{Addresses}: full CRUD with a single-default constraint enforced at the service layer; the default address auto-populates checkout.
  \item \textbf{Reviews}: one review per order (enforced); customer can update their own; 1–5 star with text; restaurant rating auto-updated.
\end{itemize}

% ╔══════════════════════════════════════════════════════╗
% ║  CHAPTER 6 — RESTAURANT PLATFORM                     ║
% ╚══════════════════════════════════════════════════════╝
\chapter{Restaurant Platform}

\section{Onboarding}

Restaurant registration is a multi-step flow:

\begin{enumerate}[leftmargin=1.5em]
  \item User registers (or logs in) and navigates to \texttt{/onboarding/restaurant}
  \item Completes the registration form including city (searched from 6,708-entry dataset)
  \item Server creates Restaurant row, sets \texttt{ownerId}, updates User's roles array, re-issues token with \texttt{restaurantId} claim
\end{enumerate}

The endpoint is idempotent: if the user already has a pending restaurant
registration, a second call returns the existing restaurant.  The auto-heal
logic on \texttt{/auth/me} re-attaches orphaned restaurantIds when token claims
drift from database state.

\section{Menu Management}

The shop menu page (commit \texttt{02c1bf8}) provides full CRUD for menu items
with drag-and-drop reordering via \texttt{@dnd-kit}.  Each item tracks:

\begin{itemize}[leftmargin=1.4em,nosep]
  \item Name, description, price (paise), category, image URL
  \item \texttt{isVeg} flag with visual indicator
  \item \texttt{isBestSeller} badge
  \item \texttt{isAvailable} toggle (inline from the menu list)
  \item \texttt{sortOrder} for custom sequence
\end{itemize}

Menu item limits are enforced via the SiteConfig \texttt{maxMenuItems} setting,
checked at creation time.

\section{Order Management}

The restaurant dashboard receives new orders via \texttt{order:new} socket events
delivered to the \texttt{shop-\{restaurantId\}} room.  Restaurant owners can
transition orders through: CONFIRMED → PREPARING → OUT\_FOR\_DELIVERY (triggers
agent assignment on the admin side).

State-machine validation in \texttt{orders.controller.js} enforces:

\begin{center}\small
\begin{tabular}{ll}
\toprule
\rowcolor{headbg}\th{Actor} & \th{Permitted Transitions} \\
\midrule
RESTAURANT & PLACED→CONFIRMED, CONFIRMED→PREPARING \\
\rowcolor{rowalt}DELIVERY   & CONFIRMED→OUT\_FOR\_DELIVERY, OUT\_FOR\_DELIVERY→DELIVERED \\
ADMIN      & Any transition, including CANCEL \\
\rowcolor{rowalt}CUSTOMER   & Any open status → CANCELLED \\
\bottomrule
\end{tabular}
\end{center}

\section{Restaurant Analytics}

The analytics endpoint returns:
\begin{itemize}[leftmargin=1.4em,nosep]
  \item Revenue by day/week/month using the \texttt{Order(restaurantId, createdAt)} composite index
  \item Order count by status
  \item Top menu items by order frequency
  \item Average order value
\end{itemize}

The Ops Intelligence Wave B scoring engine additionally computes a
restaurant performance score (0–100) based on acceptance rate, preparation
time, and review rating.

\section{Status \& Availability}

Restaurants have two availability controls:

\begin{itemize}[leftmargin=1.4em]
  \item \textbf{isOpen} (boolean): manually toggled by the owner from the shop header or settings
  \item \textbf{statusNote} (text): freeform message displayed to customers ("Back at 5pm", "Out of gas")
\end{itemize}

The \texttt{set-status} endpoint (the active one; \texttt{status} is dead code)
atomically updates both fields and emits a socket event for real-time admin
visibility.

% ╔══════════════════════════════════════════════════════╗
% ║  CHAPTER 7 — DELIVERY PLATFORM                       ║
% ╚══════════════════════════════════════════════════════╝
\chapter{Delivery Platform}

\section{Rider Onboarding}

Riders register via \texttt{POST /role/register-rider}, which adds DELIVERY to
their roles array and creates a RiderLocation row initialised to coordinates
(0, 0).  The rider then sees the delivery dashboard at \texttt{/delivery/home}.

\section{Going Online}

\texttt{PATCH /delivery/status} toggles the rider's \texttt{isAvailable} flag.
When going online, the rider's current GPS coordinates (from the frontend's
Geolocation API) are included, seeding the RiderLocation table for the admin
map.  Going online also joins the rider's Socket.IO room \texttt{agent-\{id\}},
which receives assignment events.

\section{GPS Streaming Architecture}

\begin{figure}[h!]\centering
\begin{tikzpicture}[
  box/.style={rectangle,rounded corners=4pt,draw=dark,fill=light,
    minimum width=2.5cm,minimum height=0.8cm,font=\footnotesize,align=center},
  arrow/.style={-Stealth,thick}
]
\node[box,fill=brand!15,draw=brand] (rider) at (0,0) {Rider\\Browser};
\node[box] (geoloc) at (3,0) {Geolocation\\API};
\node[box] (svc) at (6,0) {GPS Tracker\\Service};
\node[box,fill=info!15] (api) at (9,0) {POST\\\/delivery\/location};
\node[box,fill=ok!15,draw=ok] (db) at (12,0) {RiderLocation\\(Upsert)};
\node[box,fill=purple!15,draw=purple] (socket) at (9,-2) {Socket.IO\\Emit};
\node[box] (adminmap) at (5,-2) {Admin\\Live Map};
\node[box] (trackpage) at (12,-2) {Customer\\Track Page};

\draw[arrow] (rider) -- (geoloc) node[midway,above,font=\tiny] {watchPosition};
\draw[arrow] (geoloc) -- (svc) node[midway,above,font=\tiny] {coords};
\draw[arrow] (svc) -- (api) node[midway,above,font=\tiny] {POST 5s};
\draw[arrow] (api) -- (db) node[midway,above,font=\tiny] {upsert};
\draw[arrow] (api) -- (socket);
\draw[arrow,dashed] (socket) -- (adminmap) node[midway,above,font=\tiny] {rider:location:update};
\draw[arrow,dashed] (socket) -- (trackpage) node[midway,above,font=\tiny] {agent:location};
\end{tikzpicture}
\caption{GPS streaming pipeline from rider device to admin map and customer tracking page}
\end{figure}

The \texttt{riderTracker} service (extracted in commit \texttt{37593ee})
encapsulates \texttt{navigator.geolocation.watchPosition} with a 5-second
minimum interval, error handling with retry, and automatic cleanup on component
unmount.  The refactor eliminated a duplicate GPS watcher bug (commit
\texttt{81b3641}) where the same coordinates were being transmitted twice.

\section{Order Assignment}

When an admin (or the auto-assignment engine) calls \texttt{POST /admin/orders/:id/assign-delivery}:

\begin{enumerate}[leftmargin=1.5em]
  \item The system queries available DELIVERY users sorted by proximity to the restaurant
  \item Updates \texttt{Order.agentId} and \texttt{Order.estimatedDelivery}
  \item Emits \texttt{order:assigned} to the rider's \texttt{agent-\{id\}} room
  \item Emits \texttt{agent:assigned} to the order room and restaurant room
\end{enumerate}

\section{Earnings \& Sessions}

RiderSession records are created when a rider goes online and closed when they
go offline.  The earnings endpoint aggregates sessions and delivered orders to
compute:

\begin{itemize}[leftmargin=1.4em,nosep]
  \item Total online time for the period
  \item Number of deliveries completed
  \item Earnings per delivery and total earnings
  \item Session timeline for the leaderboard
\end{itemize}

% ╔══════════════════════════════════════════════════════╗
% ║  CHAPTER 8 — PAYMENTS \& FINANCE                     ║
% ╚══════════════════════════════════════════════════════╝
\chapter{Payments \& Finance}

\section{Cashfree Integration Architecture}

\begin{figure}[h!]\centering
\begin{tikzpicture}[
  box/.style={rectangle,rounded corners=4pt,draw=dark,fill=light,
    minimum width=3cm,minimum height=0.8cm,font=\footnotesize\bfseries,align=center},
  arrow/.style={-Stealth,thick},
  dashed arrow/.style={-Stealth,dashed,color=muted}
]
\node[box,fill=brand!15,draw=brand] (fe) at (0,6) {Frontend\\Checkout};
\node[box,fill=info!15,draw=info]   (be) at (5,6) {Backend\\/payments};
\node[box,fill=warn!15,draw=warn]   (cf) at (10,6) {Cashfree PG\\API};
\node[box,fill=ok!15,draw=ok]       (pg) at (5,3) {PostgreSQL\\Payment};
\node[box,fill=purple!15,draw=purple] (wh) at (10,3) {Cashfree\\Webhook};

\draw[arrow] (fe) -- (be) node[midway,above,font=\tiny] {POST create-order};
\draw[arrow] (be) -- (cf) node[midway,above,font=\tiny] {PGCreateOrder};
\draw[arrow,dashed arrow] (cf) to[bend right=15] (be) node[midway,right,font=\tiny] {session\_id};
\draw[arrow,dashed arrow] (be) to[bend right=15] (fe) node[midway,below,font=\tiny] {session\_id};
\draw[arrow] (fe) -- (cf) node[midway,above,font=\tiny,sloped] {cashfree.checkout()};
\draw[arrow] (wh) -- (be) node[midway,above,font=\tiny] {POST webhook\\+ HMAC sig};
\draw[arrow] (be) -- (pg) node[midway,left,font=\tiny] {upsert payment};
\node[font=\tiny,color=ok,right=0.1cm of wh] {HMAC-SHA256};
\end{tikzpicture}
\caption{End-to-end payment flow with Cashfree PG}
\end{figure}

\section{Webhook Security}

Cashfree delivers payment results via webhook.  The webhook handler is the most
security-critical endpoint in the platform:

\begin{lstlisting}[language=JavaScript,caption={HMAC verification (payment.controller.js)}]
// Raw body preserved before Express JSON parsing
app.use("/api/payments/webhook", express.raw({ type: "application/json" }),
  (req, _res, next) => { req.rawBody = req.body; next(); });

// In webhook handler:
const signature = req.headers["x-webhook-signature"]
               || req.headers["x-cf-signature"];
if (!signature) return res.status(400).json({ message: "Missing signature" });
// Timing-safe HMAC comparison on rawBody
\end{lstlisting}

The raw body is stored before Express JSON parsing because the HMAC is computed
over the exact wire bytes — JSON re-serialisation would produce a different hash.

\section{Refund Engine}

The refund flow supports two paths:

\begin{itemize}[leftmargin=1.4em]
  \item \textbf{Customer-initiated}: \texttt{POST /payments/my-refunds} — creates a Refund record, calls \texttt{cashfree.PGCreateRefund}, returns the Cashfree refund ID.
  \item \textbf{Admin-initiated}: \texttt{POST /payments/refunds} — same flow but without the ownership check.
\end{itemize}

A critical production bug discovered in the supply-chain audit: the service was
calling \texttt{cashfree.PGFetchRefund()} which does not exist in the SDK.  The
correct method is \texttt{PGOrderFetchRefund()}.  This meant every
\texttt{syncRefundStatus} call would 502 in production.  Fixed in commit
\texttt{00bfc3d}.

\section{Supply Chain Security}

The \texttt{cashfree-pg} npm package pinned \texttt{axios@1.15.0} — a version
with 21 reported CVEs including SSRF and prototype pollution advisories.
The audit determined that:

\begin{itemize}[leftmargin=1.4em]
  \item \textbf{SSRF}: not exploitable — Cashfree SDK calls a hardcoded API URL
  \item \textbf{Prototype pollution}: not exploitable — all user input validated by Zod before SDK
  \item \textbf{Proxy leakage}: not exploitable — no user-controlled axios configuration
\end{itemize}

Even with low exploitability, the exposure was remediated via npm \texttt{overrides}:

\begin{lstlisting}[language=json,caption={package.json axios override}]
"overrides": { "axios": "^1.17.0" }
\end{lstlisting}

This forces all transitive dependencies to use the patched version. Post-override:
\texttt{npm audit} reports 0 vulnerabilities.

\section{Payout System}

The Payouts section in the admin Finance nav (commit \texttt{a662bca}) provides
visibility into refund status and pending payouts.  Direct bank transfer
integration is a roadmap item; the current implementation tracks refund states
through the Cashfree refund API.

% ╔══════════════════════════════════════════════════════╗
% ║  CHAPTER 9 — REAL-TIME SYSTEMS                       ║
% ╚══════════════════════════════════════════════════════╝
\chapter{Real-Time Systems}

\section{Socket.IO Architecture}

Socket.IO sits alongside the Express HTTP server, sharing the Node.js \texttt{http.Server}
instance.  In production the Redis adapter (\texttt{@socket.io/redis-adapter})
is mounted to fan events across all backend instances.

\begin{lstlisting}[language=JavaScript,caption={Redis adapter initialization (socketServer.js)}]
const pub = new Redis(env.REDIS_URL);
const sub = pub.duplicate();
io.adapter(createAdapter(pub, sub));
logger.info("Socket.IO Redis adapter mounted");
\end{lstlisting}

On Redis unavailability, the adapter falls back to in-memory mode (single-instance
only), and \texttt{/health/extended} reports \texttt{adapterType: "in-memory"}.

\section{Room Design}

\begin{center}\small
\begin{tabular}{p{3.5cm}p{3cm}p{6cm}}
\toprule
\rowcolor{headbg}\th{Room} & \th{Members} & \th{Events Received} \\
\midrule
\texttt{admin}           & Admin users    & All order events, rider locations, incidents \\
\rowcolor{rowalt}\texttt{user:\{id\}}       & Logged-in user & \texttt{notification:new} \\
\texttt{shop-\{restId\}} & Restaurant owner & \texttt{order:new}, \texttt{agent:assigned} \\
\rowcolor{rowalt}\texttt{order-\{id\}}      & Customer, rider & \texttt{order:status-updated}, \texttt{agent:location} \\
\texttt{agent-\{id\}}   & Rider          & \texttt{order:assigned} \\
\rowcolor{rowalt}\texttt{delivery:\{id\}}   & Legacy alias   & (deprecated, superseded by agent-\{id\}) \\
\bottomrule
\end{tabular}
\end{center}

\section{Room Authorization — The Ownership Gate}

The critical security advance of the socket hardening sprint was converting
\texttt{order-*} rooms from capability-URL security (anyone who knows the order
ID can join) to server-side ownership verification:

\begin{lstlisting}[language=JavaScript,caption={checkOrderRoomAccess — Prisma ownership gate}]
export const checkOrderRoomAccess = async (orderId, user) => {
  if (!user) return false;
  if (!orderId || typeof orderId !== "string" || orderId.length > 100) return false;
  if (user.roles?.includes("ADMIN")) return true;
  try {
    const order = await prisma.order.findUnique({
      where: { id: orderId },
      select: { customerId: true, agentId: true, restaurantId: true },
    });
    if (!order) return false;
    if (order.customerId === user.id) return true;
    if (order.agentId === user.id) return true;
    if (user.restaurantId && order.restaurantId === user.restaurantId) return true;
    return false;
  } catch { return false; }  // fail-closed on DB error
};
\end{lstlisting}

The gate is \emph{fail-closed}: any database error denies access rather than
defaulting to permit — the only safe default for a security control.

\section{Event Map}

\begin{center}\small
\begin{tabular}{p{3.8cm}p{3.5cm}p{5.5cm}}
\toprule
\rowcolor{headbg}\th{Event} & \th{Emitter} & \th{Purpose} \\
\midrule
\texttt{order:new}            & Order service      & Notify restaurant of incoming order \\
\rowcolor{rowalt}\texttt{order:status-updated} & Order service      & Customer + admin tracking update \\
\texttt{agent:assigned}       & Admin/auto-assign  & Notify customer + restaurant of rider \\
\rowcolor{rowalt}\texttt{order:assigned}        & Order service      & Notify specific rider of their delivery \\
\texttt{rider:location:update} & Delivery route    & Admin live map position refresh \\
\rowcolor{rowalt}\texttt{agent:location}        & Delivery route     & Customer track page position update \\
\texttt{notification:new}     & Notification svc   & In-app notification badge \\
\rowcolor{rowalt}\texttt{maintenance:started}   & Config service     & All clients: show maintenance banner \\
\texttt{maintenance:ended}    & Config service     & All clients: remove banner \\
\rowcolor{rowalt}\texttt{maintenance:scheduled} & Config service     & All clients: show scheduled warning \\
\texttt{maintenance:cancelled}& Config service     & All clients: remove scheduled warning \\
\bottomrule
\end{tabular}
\end{center}

% ╔══════════════════════════════════════════════════════╗
% ║  CHAPTER 10 — SECURITY EVOLUTION                     ║
% ╚══════════════════════════════════════════════════════╝
\chapter{Security Evolution}

\section{Original Posture (Era I)}

The initial prototype had minimal security controls:

\begin{itemize}[leftmargin=1.4em]
  \item Single-role JWT with no refresh token rotation
  \item CORS wildcard (\texttt{*}) on all origins
  \item No rate limiting
  \item No input sanitisation beyond basic Express JSON parsing
  \item No HMAC verification on payment webhook
  \item No audit trail
  \item Prices trusted from client payload (price manipulation vulnerable)
\end{itemize}

The first security commit (\texttt{fef47ec}, March 28) fixed the price
manipulation vulnerability on day 2, but the other gaps persisted through Era II.

\section{Security Hardening Sprint Part 1 (June 14)}

Commit \texttt{02c7962} implemented comprehensive OWASP coverage:

\begin{center}\small
\begin{longtable}{p{4cm}p{8.5cm}}
\toprule
\rowcolor{headbg}\th{Category} & \th{Controls Implemented} \\
\midrule
Injection & Zod schema validation on all request bodies; xss library sanitises strings \\
\rowcolor{rowalt}Authentication & Helmet, CORS strict allowlist, no credentials in error messages \\
RBAC & Three-layer authorization; role checked from JWT, not request body \\
\rowcolor{rowalt}Sensitive Data & No prices/amounts from client; PII stripped from error responses in production \\
DoS Protection & express-rate-limit with Redis backing; tiered limits per endpoint sensitivity \\
\rowcolor{rowalt}Security Headers & Helmet CSP, HSTS (1yr + preload), X-Frame-Options, X-Content-Type \\
Permissions-Policy & Camera, microphone, geolocation, payment, USB — all denied \\
\rowcolor{rowalt}Info Leakage & Stack traces hidden in production; error codes not messages in API responses \\
\bottomrule
\end{longtable}
\end{center}

\section{Security Hardening Sprint Part 2 (June 14)}

Commit \texttt{0f21179} added the institutional security layer:

\begin{itemize}[leftmargin=1.4em]
  \item \textbf{AuditLog} — every admin mutation writes an immutable log entry with userId, action, resource, IP, and user-agent
  \item \textbf{Audit export} — CSV export of the audit log for compliance review
  \item \textbf{Permissions-Policy header} — restricts browser APIs the API server has no reason to use
  \item \textbf{Security regression suite} — 47 tests covering RBAC bypass, token forgery, price manipulation, IDOR, and coupon fraud
\end{itemize}

\section{IDOR Closure (June 12)}

Commit \texttt{eecf1c9} closed the IDOR gaps identified in the security audit:

\begin{tcolorbox}[crit,title=IDOR Vulnerabilities Found \& Fixed]
\begin{itemize}[leftmargin=1em,nosep]
  \item \textbf{Order IDOR}: \texttt{GET /orders/:id} returned any order to any authenticated user → Fixed with \texttt{canAccessOrder()} ownership gate
  \item \textbf{Restaurant IDOR}: \texttt{PUT /restaurants/:id} accepted updates from any RESTAURANT user → Fixed with \texttt{getRestaurantByIdAndOwner()} check
  \item \textbf{Menu IDOR}: Same pattern on all 4 menu operations → Fixed with same ownership check
  \item \textbf{Refund IDOR}: Customer could initiate refund on another customer's payment → Fixed with payment ownership verification
\end{itemize}
\end{tcolorbox}

\section{Socket Room Security (June 14)}

Before the hardening sprint, \texttt{order-*} socket rooms used capability-URL
security: any client that knew an order ID (a CUID, not truly secret) could join
and receive real-time status updates for that order.  The hardening replaced this
with a Prisma ownership query on every join, verified from the authenticated
JWT — not the client payload.

\section{JWT \& Refresh Token Security}

\begin{center}\small
\begin{tabular}{p{4cm}p{4cm}p{5cm}}
\toprule
\rowcolor{headbg}\th{Property} & \th{Before} & \th{After} \\
\midrule
Token storage      & Cookie (third-party blocked) & Authorization header in localStorage \\
\rowcolor{rowalt}Refresh token      & Stored cleartext in DB & SHA-256 hash only \\
Rotation           & None (token valid until expiry) & Rotate-on-use; replay rejected \\
\rowcolor{rowalt}Race condition     & Not handled & \texttt{deleteMany} returns count; 0 → 401 \\
Block enforcement  & Token valid until expiry & Revoke all tokens on block/suspend \\
\rowcolor{rowalt}Token payload      & userId only & userId + roles[] + activeRole + restaurantId \\
\bottomrule
\end{tabular}
\end{center}

\section{OWASP Top 10 Coverage Matrix}

\begin{center}\small
\begin{tabular}{p{4.5cm}lp{6cm}}
\toprule
\rowcolor{headbg}\th{OWASP Category} & \th{Status} & \th{Evidence} \\
\midrule
A01 Broken Access Control     & \yes{} & RBAC + IDOR gates on all resources \\
\rowcolor{rowalt}A02 Cryptographic Failures  & \yes{} & HMAC-SHA256 webhooks; bcrypt passwords; hashed tokens \\
A03 Injection                 & \yes{} & Zod validation; xss library; Prisma parameterised queries \\
\rowcolor{rowalt}A04 Insecure Design         & \yes{} & Fail-closed socket gate; server-derived prices \\
A05 Security Misconfiguration & \yes{} & Helmet; strict CORS; no stack traces in prod \\
\rowcolor{rowalt}A06 Vulnerable Components   & \yes{} & axios override; 0 audit vulnerabilities \\
A07 Auth Failures             & \yes{} & Rate limiting on /auth; token rotation; revocation \\
\rowcolor{rowalt}A08 Software Integrity      & \yes{} & npm overrides; HMAC webhook; Sentry release tags \\
A09 Logging \& Monitoring     & \yes{} & Structured JSON logs; Sentry; audit log; health endpoints \\
\rowcolor{rowalt}A10 SSRF                   & \yes{} & No user-controlled URLs in server-side requests \\
\bottomrule
\end{tabular}
\end{center}

\section{Security Test Suites}

\begin{center}\small
\begin{tabular}{lrp{7cm}}
\toprule
\rowcolor{headbg}\th{Suite} & \th{Tests} & \th{Coverage} \\
\midrule
\texttt{security.idor.test.js}        & 24 & Order/restaurant/refund IDOR gates \\
\rowcolor{rowalt}\texttt{security.regression.test.js}  & 47 & RBAC bypass, price manipulation, coupon fraud, token forgery \\
\texttt{security.sprint2.test.js}     & 18 & Input hardening, rate limiting, headers, audit trail \\
\rowcolor{rowalt}\texttt{security.platform.test.js}    & 31 & Cross-role contamination, admin endpoints, exec auth \\
\texttt{security.socket-rooms.test.js}& 14 & Socket ownership gate, edge cases, fail-closed \\
\rowcolor{rowalt}\texttt{security.supply-chain.test.js}& 19 & Webhook HMAC, SDK smoke, payment flow, refund flow \\
\midrule
\textbf{Total security tests}         & \textbf{153} & \\
\bottomrule
\end{tabular}
\end{center}

% ╔══════════════════════════════════════════════════════╗
% ║  CHAPTER 11 — OPS INTELLIGENCE & ADMIN CONSOLE      ║
% ╚══════════════════════════════════════════════════════╝
\chapter{Ops Intelligence \& Admin Console}

\section{Admin Architecture}

The admin console is a Next.js dashboard at \texttt{/admin/*} gated behind
\texttt{requireRole("ADMIN")}.  It integrates 37 API endpoints across 8
functional areas and is the operations nerve centre of the platform.

\begin{tcolorbox}[insight,title=Bootstrap Admin]
The first admin cannot be created through the normal registration flow — that
would create a chicken-and-egg problem where anyone can make themselves admin
before a real admin exists.  Commit \texttt{bc9090d} introduced \texttt{POST /admin/bootstrap-admin}:
an endpoint that is \textbf{only callable when zero admin users exist}.  After
the first admin is created, all subsequent calls return 403.
\end{tcolorbox}

\section{User Management}

\begin{center}\small
\begin{tabular}{p{4.5cm}p{8cm}}
\toprule
\rowcolor{headbg}\th{Operation} & \th{Behaviour} \\
\midrule
List users          & Paginated, searchable; shows role badges and status \\
\rowcolor{rowalt}Suspend user        & Sets \texttt{isSuspended=true}; revokes all refresh tokens immediately \\
Block user          & Sets \texttt{isBlocked=true}; prevents login; revokes all tokens \\
\rowcolor{rowalt}Unblock/Unsuspend   & Clears flag; user must re-login for new tokens \\
Delete user         & Cascades via Prisma; requires admin confirmation \\
\rowcolor{rowalt}Change role         & Appends to roles[] array; re-issues token not required (takes effect on next login) \\
\bottomrule
\end{tabular}
\end{center}

\section{Incident Engine}

Introduced in Ops Intelligence Wave A (commit \texttt{0e36a14}), the incident
engine runs every 2 minutes and detects operational anomalies:

\begin{center}\small
\begin{tabular}{p{4cm}p{3cm}p{5.5cm}}
\toprule
\rowcolor{headbg}\th{Detection Rule} & \th{Severity} & \th{Trigger Condition} \\
\midrule
Stalled order       & HIGH       & Order in PREPARING/CONFIRMED for $>$45 min \\
\rowcolor{rowalt}Payment failure spike & CRITICAL  & $>$3 failed payments in 15 minutes \\
Offline restaurant  & MEDIUM     & Restaurant closed during listed hours \\
\rowcolor{rowalt}No available riders & HIGH       & Zero riders online and open orders waiting \\
Slow delivery       & MEDIUM     & OUT\_FOR\_DELIVERY for $>$60 min \\
\bottomrule
\end{tabular}
\end{center}

Wave B (commit \texttt{d33f5cb}) added escalation:

\begin{itemize}[leftmargin=1.4em]
  \item Unacknowledged incidents auto-escalate after 30 minutes (severity upgrades one level)
  \item \texttt{escalationCount} tracks how many times an incident has escalated
  \item \texttt{IncidentEvent} provides an immutable timeline of every state change (opened, acknowledged, escalated, resolved)
\end{itemize}

\section{Performance Scoring}

Wave B's scoring endpoint computes a performance score (0–100) for each
restaurant, combining:

\begin{itemize}[leftmargin=1.4em,nosep]
  \item Acceptance rate (restaurant confirmed orders / received orders): weight 0.4
  \item Average preparation time vs. declared \texttt{estimatedDelivery}: weight 0.3
  \item Review rating (1–5 normalised to 0–100): weight 0.3
\end{itemize}

Scores power the leaderboard (top performers surfaced to customers in the
recommendations algorithm).

\section{Rider Leaderboard}

The rider leaderboard (Wave B, commit \texttt{d33f5cb}) aggregates RiderSession
data to produce:

\begin{itemize}[leftmargin=1.4em,nosep]
  \item Deliveries completed in the selected period
  \item Online hours
  \item Earnings
  \item On-time delivery rate (delivered within \texttt{estimatedDelivery})
\end{itemize}

\section{SLA Monitoring}

The \texttt{/ops/sla} endpoint returns current breach counts across:

\begin{itemize}[leftmargin=1.4em,nosep]
  \item Orders awaiting confirmation for $>$10 min
  \item Orders awaiting pickup for $>$15 min
  \item Active deliveries exceeding \texttt{estimatedDelivery}
\end{itemize}

\section{Executive Analytics}

Wave C (commit \texttt{51f0a80}) delivered the executive dashboard at
\texttt{/admin/analytics/executive} with five sub-views:

\begin{enumerate}[leftmargin=1.5em]
  \item \textbf{KPIs} — GMV, order count, average order value, active users (DAU/WAU/MAU)
  \item \textbf{Fleet} — rider count online/offline, avg delivery time, leaderboard
  \item \textbf{Customers} — new vs. returning, cohort retention, top spenders
  \item \textbf{Restaurants} — revenue ranking, acceptance rates, performance scores
  \item \textbf{Trends} — GMV by day/week, order volume heat map by hour, cuisine breakdown
\end{enumerate}

% ╔══════════════════════════════════════════════════════╗
% ║  CHAPTER 12 — SCALING ARCHITECTURE                   ║
% ╚══════════════════════════════════════════════════════╝
\chapter{Scaling Architecture}

\section{Design Premise}

GhostKitchen was designed from Era III onwards to run on multiple backend
instances without synchronisation bugs.  Two critical shared-state problems
required explicit solutions:

\begin{enumerate}[leftmargin=1.5em]
  \item \textbf{Cron jobs} — if three instances each run the incident engine independently, incidents are tripled and notifications sent three times per event.
  \item \textbf{SiteConfig cache} — if each instance caches the SiteConfig independently, a maintenance mode change takes up to 5 seconds to propagate on the instance that made the change but potentially longer on others.
\end{enumerate}

\section{Distributed Locks}

The \texttt{redisLock.js} utility implements Redis SETNX with atomic expiry:

\begin{lstlisting}[language=JavaScript,caption={Distributed lock (utils/redisLock.js)}]
export async function acquireRedisLock(key, ttlSeconds = 30) {
  const result = await redis.set(key, "locked", "EX", ttlSeconds, "NX");
  return result === "OK";
}
export async function releaseRedisLock(key) {
  await redis.del(key);
}
\end{lstlisting}

\texttt{SET NX EX} is a single atomic Redis command — there is no
TOCTOU race between checking and setting.  If the acquiring instance crashes,
the TTL ensures the lock expires automatically (housekeeping: 540s, incident:
110s, maintenance: 50s).

Three cron jobs use distributed locks:

\begin{center}\small
\begin{tabular}{p{3.5cm}p{2cm}p{3cm}p{4cm}}
\toprule
\rowcolor{headbg}\th{Job} & \th{Schedule} & \th{Lock Key} & \th{TTL} \\
\midrule
Incident engine    & \texttt{*/2 * * * *} & \texttt{lock:incident-engine}  & 110s \\
\rowcolor{rowalt}Housekeeping       & \texttt{*/10 * * * *} & \texttt{lock:housekeeping}     & 540s \\
Maintenance cron   & \texttt{* * * * *}   & \texttt{lock:maintenance-cron} & 50s \\
\bottomrule
\end{tabular}
\end{center}

\section{Three-Layer SiteConfig Cache}

\begin{figure}[h!]\centering
\begin{tikzpicture}[
  cache/.style={rectangle,rounded corners=4pt,draw=dark,fill=light,
    minimum width=3.5cm,minimum height=0.8cm,font=\footnotesize\bfseries,align=center},
  arrow/.style={-Stealth,thick}
]
\node[cache,fill=brand!15,draw=brand] (req) at (0,6) {Incoming Request};
\node[cache,fill=info!20,draw=info]   (l1)  at (0,4.5) {L1: In-Process\\(5s TTL)};
\node[cache,fill=warn!20,draw=warn]   (l2)  at (0,3)   {L2: Redis\\(30s TTL)};
\node[cache,fill=ok!15,draw=ok]       (l3)  at (0,1.5) {L3: PostgreSQL\\(source of truth)};

\draw[arrow] (req) -- (l1) node[midway,right,font=\tiny] {check};
\draw[arrow,dashed] (l1.south) -- (l2.north) node[midway,right,font=\tiny] {miss → check};
\draw[arrow,dashed] (l2.south) -- (l3.north) node[midway,right,font=\tiny] {miss → read};
\node[font=\tiny,color=muted,right=2.5cm of l1] {$<$1ms — sub-millisecond object access};
\node[font=\tiny,color=muted,right=2.5cm of l2] {$\sim$2ms — Redis network round-trip};
\node[font=\tiny,color=muted,right=2.5cm of l3] {$\sim$20ms — PostgreSQL read};
\end{tikzpicture}
\caption{Three-layer SiteConfig cache with falling-through miss chain}
\end{figure}

The cache is invalidated on write — \texttt{updateSiteConfig} clears the Redis
key and resets the in-process TTL to zero.  On the next request from any
instance, L1 and L2 are both cold and the DB read repopulates both layers.

\section{Socket.IO Horizontal Scaling}

Without the Redis adapter, a Socket.IO event emitted on instance A would only
reach clients connected to instance A.  With \texttt{@socket.io/redis-adapter},
emits are published to Redis pub/sub and all instances consume and forward them
to their locally-connected clients.

\begin{lstlisting}[language=JavaScript,caption={Redis adapter with dual-client ioredis}]
const pub = new Redis(env.REDIS_URL);
const sub = pub.duplicate();   // separate connection for subscribe
io.adapter(createAdapter(pub, sub));
\end{lstlisting}

Two clients are required: Redis does not allow a connection to both publish and
subscribe simultaneously.  The \texttt{duplicate()} call inherits all connection
options from the pub client.

\section{Upstash REST Fallback}

The Redis client module supports two connection modes, selected by environment
variable:

\begin{itemize}[leftmargin=1.4em]
  \item \textbf{\texttt{REDIS\_URL} set}: ioredis TCP connection (used for Socket.IO adapter; pub/sub required)
  \item \textbf{\texttt{UPSTASH\_REDIS\_REST\_URL} set}: Upstash REST client (HTTP-only hosting environments; no pub/sub)
\end{itemize}

The fallback degrades gracefully: if both are unavailable, the platform boots in
single-instance mode with in-memory socket state, logging a WARNING rather than
crashing.

\section{Horizontal Scale Checklist}

\begin{center}\small
\begin{tabular}{p{5cm}lp{6cm}}
\toprule
\rowcolor{headbg}\th{Concern} & \th{Status} & \th{Solution} \\
\midrule
Cron duplication          & \yes{} & Distributed Redis locks \\
\rowcolor{rowalt}Socket fan-out            & \yes{} & Redis adapter pub/sub \\
SiteConfig consistency    & \yes{} & Cache invalidation on write; 30s max lag \\
\rowcolor{rowalt}Session sharing           & \yes{} & JWTs (stateless); refresh tokens in DB \\
Rate limiting             & \half{} & In-process; Redis-backed RL planned \\
\rowcolor{rowalt}File uploads              & \half{} & Not implemented (image URLs only) \\
Sticky sessions           & \yes{} & Not required (stateless design) \\
\bottomrule
\end{tabular}
\end{center}

% ╔══════════════════════════════════════════════════════╗
% ║  CHAPTER 13 — OBSERVABILITY                          ║
% ╚══════════════════════════════════════════════════════╝
\chapter{Observability}

\section{The Three Pillars}

GhostKitchen's observability stack implements all three pillars:
logs (Winston), traces (X-Request-ID correlation), and errors (Sentry).

\section{Structured Logging}

The Winston logger emits in two formats:

\begin{itemize}[leftmargin=1.4em]
  \item \textbf{Production} (JSON): machine-parseable, compatible with log aggregators (Datadog, Papertrail, Render Logs)
  \item \textbf{Development} (printf): coloured, human-readable with timestamp prefix
\end{itemize}

\begin{lstlisting}[language=JavaScript,caption={Winston configuration (utils/logger.js)}]
const consoleFormat = isProduction
  ? winston.format.combine(
      winston.format.errors({ stack: true }),
      winston.format.json()        // pure JSON — no printf
    )
  : winston.format.combine(
      winston.format.colorize(),
      winston.format.timestamp({ format: "HH:mm:ss" }),
      winston.format.printf(({ level, message, timestamp, ...meta }) =>
        `${timestamp} ${level}: ${message} ${
          Object.keys(meta).length ? JSON.stringify(meta) : ""
        }`
      )
    );
\end{lstlisting}

\section{Request Tracing}

The request tracing middleware assigns a UUID to every incoming request:

\begin{enumerate}[leftmargin=1.5em]
  \item Reads \texttt{X-Request-ID} from incoming headers (if the gateway already assigned one) or generates \texttt{crypto.randomUUID()}
  \item Attaches \texttt{req.id} for downstream use
  \item Sets \texttt{X-Request-ID} response header for client-side correlation
  \item On \texttt{res.finish}: logs a structured access record with requestId, method, path, statusCode, durationMs, ip, userAgent, userId
\end{enumerate}

Every Sentry capture includes \texttt{requestId} as a tag, enabling a single
search in Sentry to retrieve the exact Winston log line that preceded an error.

\section{Sentry Integration}

\begin{center}\small
\begin{tabular}{p{4cm}p{8.5cm}}
\toprule
\rowcolor{headbg}\th{Feature} & \th{Implementation} \\
\midrule
Release tracking  & \texttt{release: process.env.SENTRY\_RELEASE || process.env.RENDER\_GIT\_COMMIT || "unknown"} \\
\rowcolor{rowalt}Error capture     & \texttt{captureException(err, \{ requestId, userId \})} from all error surfaces \\
Request ID tag    & \texttt{scope.setTag("requestId", requestId)} — correlates Sentry with logs \\
\rowcolor{rowalt}User context      & \texttt{scope.setUser(\{ id: userId \})} — privacy-safe (no email in Sentry) \\
Cron visibility   & Incident engine, housekeeping, and maintenance-cron all report to Sentry \\
\rowcolor{rowalt}Global safety net & \texttt{process.on("unhandledRejection")} and \texttt{"uncaughtException"} captured \\
\bottomrule
\end{tabular}
\end{center}

\section{Health Endpoints}

Two health endpoints are registered before all middleware (so they always respond):

\subsection{\texttt{GET /health} — Readiness Check}

Returns 200 OK (or 503) based on DB and Redis connectivity:

\begin{lstlisting}[language=json,caption={/health response when degraded}]
{
  "status": "DEGRADED",
  "timestamp": "2026-06-14T12:00:00.000Z",
  "release": "a7b4c9d",
  "checks": {
    "database": { "status": "UP", "responseTimeMs": 12 },
    "redis":    { "status": "DOWN", "error": "Connection refused" }
  }
}
\end{lstlisting}

Status semantics: \texttt{DOWN} (503) if DB fails; \texttt{DEGRADED} (200) if
Redis fails; \texttt{OK} (200) if both healthy.

\subsection{\texttt{GET /health/extended} — Deep Diagnostic}

Returns everything the readiness check reports plus:

\begin{itemize}[leftmargin=1.4em,nosep]
  \item Maintenance mode state (on/off, reason, scheduledAt, endsAt)
  \item Socket adapter type (\texttt{"redis"} or \texttt{"in-memory"})
  \item Cron job status: lastRunAt, lastError for incident engine and housekeeping
  \item Release version
\end{itemize}

This endpoint is intended for on-call engineers — it exposes no authentication
tokens or user data, only operational metadata.

\section{Observability Gap Closure}

\begin{center}\small
\begin{tabular}{p{4cm}p{4.5cm}p{4cm}}
\toprule
\rowcolor{headbg}\th{Gap} & \th{Before Hardening} & \th{After Hardening} \\
\midrule
Log format            & Colourised printf in prod  & Pure JSON in prod \\
\rowcolor{rowalt}Request correlation   & No request IDs             & UUID on every request \\
Sentry release        & "unknown" always           & Git commit SHA \\
\rowcolor{rowalt}Cron visibility       & Silent on error            & Sentry capture + status export \\
Health depth          & \{status: "ok"\} only      & Full diagnostic payload \\
\rowcolor{rowalt}Error context         & Error only                 & Error + requestId + userId \\
Socket adapter vis.   & Not exposed                & adapterType in /health/extended \\
\bottomrule
\end{tabular}
\end{center}

% ╔══════════════════════════════════════════════════════╗
% ║  CHAPTER 14 — TESTING ARCHITECTURE                   ║
% ╚══════════════════════════════════════════════════════╝
\chapter{Testing Architecture}

\section{Test Strategy}

GhostKitchen uses Vitest as the test runner on both backend and frontend, with
Supertest for HTTP integration tests and the React Testing Library for
component tests.

\begin{center}\small
\begin{tabular}{lrrrr}
\toprule
\rowcolor{headbg}\th{Suite} & \th{Files} & \th{Tests} & \th{Pass Rate} & \th{Coverage} \\
\midrule
Backend (unit + integration) & 36 & 927 & 100\% & 76.6\% \\
\rowcolor{rowalt}Frontend (component + E2E)  & 6  & 109 & 100\% & Partial \\
\textbf{Total}               & \textbf{42} & \textbf{1,036} & \textbf{100\%} & — \\
\bottomrule
\end{tabular}
\end{center}

\section{Coverage Journey}

The backend test coverage story is one of the most dramatic arcs in the project:

\begin{center}\small
\begin{tabular}{lrl}
\toprule
\rowcolor{headbg}\th{Date} & \th{Coverage} & \th{What Changed} \\
\midrule
June 3  & 0\%    & Rebuild started; no tests inherited from Era I \\
\rowcolor{rowalt}June 11 & 17.6\% & First test files committed (\texttt{d8c97e4}) \\
June 12 & 24.8\% & 103 tests added across payment, auth, orders \\
\rowcolor{rowalt}June 13 & 76.6\% & 927-test marathon sprint (commits \texttt{ed0b03d}–\texttt{cd7462b}) \\
\bottomrule
\end{tabular}
\end{center}

The 76.6\% figure represents 36 test files and 927 assertions landing in a
single day — a rate of approximately one test every 10 minutes of actual
development time.

\section{Backend Test Files}

\begin{center}\footnotesize
\begin{longtable}{p{5.5cm}p{7cm}}
\toprule
\rowcolor{headbg}\th{File} & \th{Coverage} \\
\midrule
\texttt{auth.test.js}                  & Registration, login, token refresh, logout, concurrent rotation \\
\rowcolor{rowalt}\texttt{orders.test.js}                & Placement, status transitions, IDOR, reorder \\
\texttt{payment.test.js}               & Create, verify, webhook idempotency, refund \\
\rowcolor{rowalt}\texttt{restaurant.test.js}            & CRUD, menu management, availability, analytics \\
\texttt{delivery.test.js}              & Online/offline, location ping, assignment, earnings \\
\rowcolor{rowalt}\texttt{admin.test.js}                 & User management, bootstrap admin, operations \\
\texttt{security.idor.test.js}         & 24 IDOR scenarios across all resource types \\
\rowcolor{rowalt}\texttt{security.regression.test.js}   & 47 regression tests for known vulnerability patterns \\
\texttt{security.sprint2.test.js}      & 18 hardening tests (headers, rate limiting, audit) \\
\rowcolor{rowalt}\texttt{security.platform.test.js}     & 31 cross-role contamination tests \\
\texttt{security.socket-rooms.test.js} & 14 socket ownership gate tests \\
\rowcolor{rowalt}\texttt{security.supply-chain.test.js} & 19 supply chain and webhook security tests \\
\texttt{middleware.extended.test.js}   & Tracing, error handler, CORS, sanitisation \\
\rowcolor{rowalt}\texttt{socket.rooms.test.js}          & Room join/leave, event fan-out, Redis adapter \\
\texttt{incident.job.test.js}          & Detection rules, escalation, Sentry integration \\
\rowcolor{rowalt}\texttt{housekeeping.job.test.js}       & Payment expiry, token purge, distributed lock \\
\texttt{config.service.test.js}        & Three-layer cache, maintenance mode, socket events \\
\rowcolor{rowalt}\texttt{ops.test.js}                   & Performance scoring, leaderboard, SLA monitoring \\
\bottomrule
\end{longtable}
\end{center}

\section{CI/CD Pipeline}

The GitHub Actions pipeline (commit \texttt{f065ac6}) runs on every push:

\begin{lstlisting}[language=yaml,caption={CI pipeline (.github/workflows/ci.yml)}]
jobs:
  backend-tests:
    runs-on: ubuntu-latest
    services:
      postgres: { image: postgres:15, env: { POSTGRES_DB: ghostkitchen_test } }
      redis:    { image: redis:7 }
    steps:
      - uses: actions/checkout@v4
      - uses: actions/setup-node@v4 with: node-version: "20"
      - run: npm ci
      - run: npx prisma migrate deploy
      - run: npm test -- --coverage
      - uses: codecov/codecov-action@v4
  frontend-tests:
    runs-on: ubuntu-latest
    steps:
      - uses: actions/checkout@v4
      - run: npm ci && npm test
\end{lstlisting}

The pipeline tests against real PostgreSQL and Redis services (not mocks),
following the lesson from Era II where mocked tests missed production failures.

\section{Vitest Configuration}

\begin{lstlisting}[language=JavaScript,caption={vitest.config.js}]
export default defineConfig({
  test: {
    environment: "node",
    globals: true,
    setupFiles: ["./src/__tests__/setup.js"],
    coverage: {
      provider: "v8",
      reporter: ["text", "json", "html"],
      exclude: ["node_modules", "coverage", "src/__tests__"],
      thresholds: { lines: 70, functions: 70 },
    },
    poolOptions: { threads: { singleThread: true } },
  },
});
\end{lstlisting}

\texttt{singleThread: true} prevents test isolation failures from parallel
database access across test files.

% ╔══════════════════════════════════════════════════════╗
% ║  CHAPTER 15 — MAJOR BUGS & RESOLUTIONS               ║
% ╚══════════════════════════════════════════════════════╝
\chapter{Major Bugs \& Resolutions}

\section{Hall of Fame: The Worst Bugs}

This chapter documents the most significant bugs encountered during GhostKitchen's
development — not to catalogue mistakes, but because each one produced a durable
engineering lesson that influenced the final architecture.

\subsection{Bug 1: The NextAuth Cookie Crisis (April 7)}

\textbf{Symptom}: All API calls returned 401 in production despite the user being
logged in.  Worked perfectly on localhost.

\textbf{Root Cause}: NextAuth stored session data in cookies.  Vercel (frontend)
and Render (backend) operate on different domains.  Safari and Firefox block
third-party cookies by default.  The session cookie set by Vercel was never
sent to Render.

\textbf{Resolution}: Replaced NextAuth with custom JWT tokens passed via
Authorization header.  Tokens stored in localStorage.  This is intentional and
documented — it bypasses the Safari restriction because Authorization headers
are not subject to cookie third-party blocking.

\textbf{Lesson}: Cross-origin deployments cannot rely on cookies for authentication
unless both origins share a parent domain.

\subsection{Bug 2: The Payment Race Condition (April 7)}

\textbf{Symptom}: Duplicate payments created for single orders.  Cashfree's
webhook retry logic was sending the same webhook 2–3 times within seconds.

\textbf{Root Cause}: The webhook handler used SELECT-then-UPDATE: check if the
payment exists, then create or update it.  Two concurrent webhook requests both
passed the SELECT (neither payment existed yet) and both proceeded to INSERT.

\textbf{Resolution}: Introduced the \texttt{PaymentWebhook} idempotency table with
a unique key \texttt{eventId = \$\{orderId\}\_\$\{status\}\_\$\{cfPaymentId\}}.
The first request inserts; any duplicate \texttt{findUnique} finds the existing
record and returns early.

\textbf{Lesson}: Webhook handlers must be idempotent.  External payment providers
will retry on timeout or failure; the handler must handle duplicates.

\subsection{Bug 3: The Token Loop (June 4)}

\textbf{Symptom}: Logging in with the new JWT system on the frontend caused an
infinite re-render loop.  The login page kept flashing.

\textbf{Root Cause}: The Zustand auth store had a hydration timing issue.  On
mount, it would restore tokens from localStorage, then immediately call
\texttt{/auth/me} to validate — but if this call returned 401 for any reason
(including a legitimate expiry), the store would clear state, which triggered
the component to redirect to login, which triggered another mount, repeat.

\textbf{Resolution}: Added a \texttt{hydrated} boolean to the Zustand store.
All auth redirect logic was gated behind \texttt{if (!hydrated) return null}.
Six commits across June 4–5 progressively tightened the hydration guard.

\textbf{Lesson}: Persisted client-side state that depends on server validation
requires a stable hydration gate before auth redirects fire.

\subsection{Bug 4: The Ghost Restaurant (June 5)}

\textbf{Symptom}: Restaurant owners who registered their restaurant would lose
their restaurant dashboard after the next token refresh.  The token was correctly
issued at registration, but the refreshed token had no \texttt{restaurantId}.

\textbf{Root Cause}: The token refresh service re-derived the token payload from
\texttt{buildTokenPayload(user)} — which pulled from the User record.  The
User's \texttt{restaurantId} field was only updated on the Restaurant row, not
the User row.

\textbf{Resolution}: The \texttt{/auth/me} endpoint now auto-heals: if the user
has RESTAURANT in their roles array but no \texttt{restaurantId} in their
decoded token, the server queries \texttt{Restaurant.findFirst(\{where:\{ownerId\}\})}
and re-issues a corrected token.

\textbf{Lesson}: Derived state in JWTs becomes stale.  Any field that can change
after token issuance needs a re-validation path.

\subsection{Bug 5: The OOM Crash (June 12)}

\textbf{Symptom}: Admin analytics endpoint caused the Render server to run out of
memory on large datasets and restart.

\textbf{Root Cause}: The analytics query fetched all Order rows for a restaurant
with no date range filter — on a restaurant with thousands of orders, this
loaded megabytes of data into Node.js heap for in-process aggregation.

\textbf{Resolution}: Date range filters (\texttt{startDate/endDate}) are now
required parameters on analytics endpoints.  The query pushes aggregation into
PostgreSQL (using Prisma's \texttt{groupBy} and \texttt{\_sum}) rather than
loading raw rows.

\textbf{Lesson}: Never load unbounded result sets into application memory.
Push aggregation to the database; paginate everything else.

\subsection{Bug 6: The Refund SDK Miss (June 12)}

\textbf{Symptom}: Every refund status sync call returned 502 in production.

\textbf{Root Cause}: The refund sync service called \texttt{cashfree.PGFetchRefund()},
which does not exist in the \texttt{cashfree-pg} SDK v4.  The correct method
is \texttt{cashfree.PGOrderFetchRefund()}.  In development, the code path was
never hit (no real Cashfree credentials); the bug was invisible until the
production refund flow was exercised.

\textbf{Resolution}: Corrected the method call and added a test that stubs the
SDK and verifies the correct method name is called.

\textbf{Lesson}: SDK method names are not stable across major versions.  Any call
to an external SDK should be covered by a test that verifies the method exists
and is called with the expected signature.

\subsection{Bug 7: The Duplicate GPS Watcher (June 14)}

\textbf{Symptom}: Rider GPS position updates were arriving at the server twice for
every physical movement, causing double notifications on the admin map.

\textbf{Root Cause}: The delivery dashboard had two separate \texttt{useEffect}
hooks, each calling \texttt{navigator.geolocation.watchPosition}.  Both watchers
triggered on every GPS update.

\textbf{Resolution}: Commit \texttt{81b3641} extracted all GPS tracking logic
into a single \texttt{riderTracker} service module with a singleton watcher.
The \texttt{useEffect} cleanup function calls \texttt{watchPosition.clear()} to
prevent stale watchers on re-mount.

\textbf{Lesson}: \texttt{watchPosition} must be cleaned up on unmount.  Browser
geolocation watchers persist across React re-renders unless explicitly cleared.

% ╔══════════════════════════════════════════════════════╗
% ║  CHAPTER 16 — FINAL ARCHITECTURE                     ║
% ╚══════════════════════════════════════════════════════╝
\chapter{Final Architecture}

\section{System Topology}

\begin{figure}[h!]\centering
\begin{tikzpicture}[
  svc/.style={rectangle,rounded corners=5pt,draw=dark,fill=light,
    minimum width=3.5cm,minimum height=1cm,font=\small\bfseries,align=center},
  ext/.style={rectangle,rounded corners=5pt,draw=muted,fill=muted!10,
    minimum width=3.5cm,minimum height=1cm,font=\small,align=center},
  arrow/.style={-Stealth,thick},
  darrow/.style={Stealth-Stealth,thick,dashed,color=muted}
]
% User tier
\node[svc,fill=info!15,draw=info] (cust) at (0,9)   {Customer\\Browser};
\node[svc,fill=info!15,draw=info] (rest) at (4,9)   {Restaurant\\Browser};
\node[svc,fill=info!15,draw=info] (rider) at (8,9)  {Rider\\Mobile};
\node[svc,fill=info!15,draw=info] (admin) at (12,9) {Admin\\Browser};

% CDN
\node[ext] (vercel) at (6,7) {Vercel Edge\\(Next.js 15 SSR)};

% API
\node[svc,fill=brand!15,draw=brand] (render) at (6,5) {Render\\(Node.js 20 Express)};

% Storage
\node[svc,fill=ok!15,draw=ok]       (pg)    at (2,3)  {Neon PostgreSQL\\(Primary)};
\node[svc,fill=warn!15,draw=warn]   (redis) at (6,3)  {Redis / Upstash\\(Cache + PubSub)};
\node[svc,fill=purple!15,draw=purple] (sentry) at (10,3) {Sentry\\(Error Tracking)};

% External
\node[ext] (cf)    at (2,1)  {Cashfree PG\\(Payments)};
\node[ext] (resend) at (6,1) {Resend\\(Email)};
\node[ext] (render2) at (10,1) {Render Logs\\(Structured JSON)};

% Arrows
\draw[arrow] (cust) -- (vercel);
\draw[arrow] (rest) -- (vercel);
\draw[arrow] (rider) -- (vercel);
\draw[arrow] (admin) -- (vercel);
\draw[arrow] (vercel) -- (render) node[midway,right,font=\tiny] {REST + WS};
\draw[arrow] (render) -- (pg);
\draw[arrow] (render) -- (redis);
\draw[arrow] (render) -- (sentry);
\draw[arrow] (render) -- (cf);
\draw[arrow] (render) -- (resend);
\draw[arrow,dashed,color=muted] (render) -- (render2);
\draw[darrow] (cf) to[bend left=20] (render) node[midway,left,font=\tiny] {webhook};
\end{tikzpicture}
\caption{Complete production system topology}
\end{figure}

\section{Dependency Inventory}

\subsection{Backend Dependencies (23 runtime)}

\begin{center}\footnotesize
\begin{tabular}{p{4.5cm}p{3.5cm}p{4.5cm}}
\toprule
\rowcolor{headbg}\th{Package} & \th{Version} & \th{Purpose} \\
\midrule
express               & \^{}4.18         & HTTP framework \\
\rowcolor{rowalt}@prisma/client        & \^{}6.9          & Database ORM \\
jsonwebtoken          & \^{}9.0          & JWT signing/verification \\
\rowcolor{rowalt}bcryptjs              & \^{}2.4          & Password hashing \\
socket.io             & \^{}4.8          & WebSocket server \\
\rowcolor{rowalt}@socket.io/redis-adapter & \^{}8.3      & Multi-instance socket fan-out \\
ioredis               & \^{}5.6          & Redis client (TCP) \\
\rowcolor{rowalt}@upstash/redis        & \^{}1.34         & Redis REST client (fallback) \\
node-cron             & \^{}3.0          & Cron job scheduler \\
\rowcolor{rowalt}cashfree-pg           & \^{}4.0          & Cashfree payment SDK \\
@sentry/node          & \^{}8.52         & Error tracking \\
\rowcolor{rowalt}winston               & \^{}3.17         & Structured logging \\
helmet                & \^{}8.0          & HTTP security headers \\
\rowcolor{rowalt}cors                  & \^{}2.8          & CORS middleware \\
express-rate-limit    & \^{}7.5          & Rate limiting \\
\rowcolor{rowalt}compression           & \^{}1.7          & gzip responses \\
zod                   & \^{}3.23         & Request validation schemas \\
\rowcolor{rowalt}xss                   & \^{}1.0          & Input sanitisation \\
resend                & \^{}4.5          & Transactional email \\
\rowcolor{rowalt}uuid                  & \^{}11.0         & Request ID generation \\
dotenv                & \^{}16.5         & Env var loading \\
\rowcolor{rowalt}axios                 & \^{}1.17         & HTTP client (overridden) \\
multer                & \^{}1.4          & Multipart form handling \\
\bottomrule
\end{tabular}
\end{center}

\subsection{Frontend Dependencies (26 runtime)}

\begin{center}\footnotesize
\begin{tabular}{p{4.5cm}p{3cm}p{5cm}}
\toprule
\rowcolor{headbg}\th{Package} & \th{Version} & \th{Purpose} \\
\midrule
next                  & 15.2.4       & App router framework \\
\rowcolor{rowalt}react                 & \^{}19.0        & UI rendering \\
zustand               & \^{}5.0         & Client state management \\
\rowcolor{rowalt}@tanstack/react-query & \^{}5.75        & Server state, caching \\
socket.io-client      & \^{}4.8         & WebSocket client \\
\rowcolor{rowalt}axios                 & \^{}1.8         & HTTP client \\
tailwindcss           & \^{}3.4         & Utility CSS \\
\rowcolor{rowalt}leaflet               & \^{}1.9         & Map rendering \\
react-leaflet         & \^{}4.2         & React Leaflet bindings \\
\rowcolor{rowalt}@dnd-kit/core         & \^{}6.3         & Drag-and-drop \\
@sentry/nextjs        & \^{}8.52        & Frontend error tracking \\
\rowcolor{rowalt}react-hot-toast       & \^{}2.5         & Toast notifications \\
lucide-react          & \^{}0.511       & Icon library \\
\rowcolor{rowalt}date-fns              & \^{}4.1         & Date formatting \\
recharts              & \^{}2.15        & Analytics charts \\
\rowcolor{rowalt}clsx                  & \^{}2.1         & Conditional classnames \\
\bottomrule
\end{tabular}
\end{center}

\section{Environment Configuration}

\begin{center}\small
\begin{tabular}{p{5cm}p{7.5cm}}
\toprule
\rowcolor{headbg}\th{Variable} & \th{Purpose} \\
\midrule
\texttt{DATABASE\_URL}          & Prisma connection string (Neon/Supabase) \\
\rowcolor{rowalt}\texttt{JWT\_SECRET}            & Access token signing key \\
\texttt{REFRESH\_TOKEN\_SECRET} & Refresh token signing key (separate secret) \\
\rowcolor{rowalt}\texttt{REDIS\_URL}             & ioredis connection string \\
\texttt{UPSTASH\_REDIS\_REST\_URL} & Upstash REST endpoint (fallback) \\
\rowcolor{rowalt}\texttt{UPSTASH\_REDIS\_REST\_TOKEN} & Upstash authentication token \\
\texttt{CASHFREE\_CLIENT\_ID}   & Cashfree API key \\
\rowcolor{rowalt}\texttt{CASHFREE\_CLIENT\_SECRET} & Cashfree API secret \\
\texttt{CASHFREE\_ENV}          & \texttt{PRODUCTION} or \texttt{SANDBOX} \\
\rowcolor{rowalt}\texttt{CASHFREE\_WEBHOOK\_SECRET} & HMAC-SHA256 webhook signing key \\
\texttt{RESEND\_API\_KEY}       & Transactional email API key \\
\rowcolor{rowalt}\texttt{SENTRY\_DSN}            & Sentry project endpoint \\
\texttt{SENTRY\_RELEASE}        & Build identifier (commit SHA) \\
\rowcolor{rowalt}\texttt{RENDER\_GIT\_COMMIT}    & Render-injected commit SHA (fallback) \\
\texttt{CLIENT\_URL}            & Frontend origin for CORS allowlist \\
\rowcolor{rowalt}\texttt{NODE\_ENV}              & \texttt{production} | \texttt{development} | \texttt{test} \\
\texttt{PORT}                   & HTTP listen port (default 10000) \\
\rowcolor{rowalt}\texttt{FRONTEND\_URL}          & Frontend URL for email links \\
\bottomrule
\end{tabular}
\end{center}

% ╔══════════════════════════════════════════════════════╗
% ║  CHAPTER 17 — PRODUCTION READINESS                   ║
% ╚══════════════════════════════════════════════════════╝
\chapter{Production Readiness}

\section{Deployment Pipeline}

\subsection{Backend — Render}

Render detects pushes to \texttt{main} and runs the build command:

\begin{lstlisting}[language=bash]
# Build command
npm install && npx prisma generate && npx prisma migrate deploy

# Start command
node src/server.js
\end{lstlisting}

\texttt{prisma migrate deploy} applies any pending migrations without
interactive prompts and without resetting data — safe for production.

\subsection{Frontend — Vercel}

Vercel's zero-config Next.js detection handles the frontend.  The only
configuration is the environment variable for the backend API URL
(\texttt{NEXT\_PUBLIC\_API\_URL}).

\section{Self-Healing Migrations}

Render can lock the deployment process if a previous migration deployment
crashed mid-run, leaving Prisma's \texttt{\_prisma\_migrations} table in an
error state.  This manifests as error code \texttt{P3009}.

The repair script (\texttt{scripts/repair-migration.js}) executes:

\begin{lstlisting}[language=JavaScript,caption={P3009 repair script}]
await prisma.$executeRaw`
  DELETE FROM "_prisma_migrations"
  WHERE "finished_at" IS NULL
`;
\end{lstlisting}

This removes the stuck migration record and allows \texttt{prisma migrate deploy}
to re-run.  The script is documented in the README and linked from the Render
deployment docs.

\section{Maintenance Mode}

The SiteConfig \texttt{maintenanceMode} flag causes the global middleware to
return 503 on all routes except \texttt{/health} and \texttt{/auth/login}.
This enables zero-downtime maintenance without stopping the server process.

\begin{lstlisting}[language=JavaScript,caption={Maintenance mode gate}]
if (config.maintenanceMode && !ALLOWED_DURING_MAINTENANCE.includes(req.path)) {
  return res.status(503).json({
    message: config.maintenanceReason || "Platform under maintenance",
    retry_after: 300,
  });
}
\end{lstlisting}

Scheduled maintenance uses \texttt{maintenanceScheduledAt} and
\texttt{maintenanceEndsAt}: a frontend banner appears 24h before the window,
and the cron job atomically flips \texttt{maintenanceMode} at the scheduled time.

\section{Production Readiness Checklist}

\begin{center}\small
\begin{tabular}{p{5.5cm}lp{6cm}}
\toprule
\rowcolor{headbg}\th{Criterion} & \th{Status} & \th{Evidence} \\
\midrule
\multicolumn{3}{l}{\textbf{Security}} \\
\midrule
OWASP Top 10         & \yes{} & Full coverage matrix (Chapter 10) \\
\rowcolor{rowalt}SQL injection         & \yes{} & Prisma parameterised queries only \\
Auth hardening        & \yes{} & Token rotation, revocation, bcrypt \\
\rowcolor{rowalt}Input validation      & \yes{} & Zod + xss on all endpoints \\
Rate limiting         & \yes{} & Tiered per endpoint sensitivity \\
\rowcolor{rowalt}CORS lockdown         & \yes{} & Strict allowlist + preview pattern \\
Dependency audit      & \yes{} & 0 npm audit findings \\
\midrule
\multicolumn{3}{l}{\textbf{Reliability}} \\
\midrule
Health endpoints      & \yes{} & \texttt{/health} + \texttt{/health/extended} \\
\rowcolor{rowalt}Distributed locks     & \yes{} & All 3 cron jobs locked \\
Graceful shutdown     & \yes{} & SIGTERM handler + DB disconnect \\
\rowcolor{rowalt}Error capture         & \yes{} & Sentry + unhandledRejection + uncaughtException \\
Maintenance mode      & \yes{} & Scheduled + manual + auto-end \\
\midrule
\multicolumn{3}{l}{\textbf{Observability}} \\
\midrule
Structured logging    & \yes{} & Winston JSON in production \\
\rowcolor{rowalt}Request tracing       & \yes{} & X-Request-ID on every request \\
Sentry release tags   & \yes{} & Git SHA attached to every event \\
\rowcolor{rowalt}Cron visibility       & \yes{} & lastRunAt, lastError exported \\
\midrule
\multicolumn{3}{l}{\textbf{Data Integrity}} \\
\midrule
Migration strategy    & \yes{} & \texttt{migrate deploy} (never db push) \\
\rowcolor{rowalt}Idempotent DDL        & \yes{} & IF NOT EXISTS from migration 5 \\
Payment idempotency   & \yes{} & Webhook deduplication table \\
\rowcolor{rowalt}Coupon atomicity      & \yes{} & Prisma transaction on order + coupon decrement \\
\midrule
\multicolumn{3}{l}{\textbf{Testing}} \\
\midrule
Test coverage         & \yes{} & 76.6\% backend, 1,036 total tests \\
\rowcolor{rowalt}Security test suite   & \yes{} & 153 tests across 6 security files \\
CI/CD pipeline        & \yes{} & GitHub Actions on every push \\
\rowcolor{rowalt}Integration tests     & \yes{} & Real PostgreSQL + Redis in CI \\
\bottomrule
\end{tabular}
\end{center}

\section{Outstanding Items}

\begin{tcolorbox}[warn,title=Known Gaps (Non-Blocking)]
\begin{enumerate}[leftmargin=1.4em,nosep]
  \item \textbf{Frontend test coverage}: 109 tests covering 6 of 44 pages; 38 pages untested
  \item \textbf{Redis-backed rate limiting}: current in-process rate limiting resets on deploy
  \item \textbf{Image upload}: image URLs only; no CDN or file storage integration
  \item \textbf{Dead code}: 5 minor react-hooks lint warnings; 6 dead-code items documented in audit
  \item \textbf{Payment webhook retry}: no exponential backoff or dead-letter queue for failed webhook processing
  \item \textbf{Database backups}: Neon provides point-in-time recovery but no explicit backup schedule documented
\end{enumerate}
\end{tcolorbox}

% ╔══════════════════════════════════════════════════════╗
% ║  CHAPTER 18 — LESSONS LEARNED                        ║
% ╚══════════════════════════════════════════════════════╝
\chapter{Lessons Learned}

\section{On Architecture}

\subsection{Cross-origin authentication needs explicit design}
NextAuth's cookie-based session management is excellent within a monorepo
deployment.  When the frontend and backend run on different domains — as they do
in the Vercel + Render deployment model — cookies become unreliable.  Design the
auth architecture for the deployment topology from day one.

\subsection{JSON Web Tokens are a cache, not a source of truth}
The ``Ghost Restaurant'' bug (Chapter 15) demonstrated that any state encoded in
a JWT becomes stale the moment the underlying data changes.  A JWT's claims
should be treated as a snapshot that might be up to 15 minutes old.  Critical
state (\texttt{isBlocked}, active restaurant) must be validated against the
database at token-refresh time.

\subsection{In-memory state is invisible in multi-instance deployments}
Both the cron duplication bug and the socket fan-out problem had the same root
cause: in-memory state that was correct for single-instance development but
wrong for multi-instance production.  Redis is not a performance optimisation
in this context — it is a correctness requirement.

\section{On Database Design}

\subsection{Add idempotency guards from migration 1, not migration 15}
The self-healing \texttt{IF NOT EXISTS} DDL should have been the default from
the first migration.  Production databases cannot be reset; every migration must
be safe to replay.

\subsection{Push aggregation into the database}
The OOM bug (Chapter 15) was caused by loading raw rows for in-process
aggregation.  PostgreSQL can aggregate millions of rows in milliseconds.  Never
load a result set into application memory when the database can compute the
answer.

\subsection{Index for the queries you actually run}
The order analytics index on \texttt{(restaurantId, createdAt)} improved
analytics query performance by 10–50× on large datasets.  Indexing should be
driven by query plans, not intuition.

\section{On Security}

\subsection{IDOR is the highest-probability vulnerability class}
Of all the security gaps closed, IDOR (Insecure Direct Object Reference) had
the highest ratio of instances to time-to-find.  Every resource endpoint that
accepts an ID parameter should have an ownership check, always.

\subsection{Webhook security is non-negotiable}
The payment webhook is the most security-critical endpoint in a payment-enabled
application.  HMAC verification, raw body preservation, and idempotency are not
optional.  They should be in place before the first line of business logic.

\subsection{Security controls must be fail-closed}
The socket room ownership gate returns \texttt{false} on any error, including
database errors.  This is the only safe default.  A security control that
``fails open'' on error provides no protection in the scenarios where something
has already gone wrong.

\section{On Testing}

\subsection{Test what you deploy, not what you mock}
The CI pipeline runs against real PostgreSQL and Redis — not mocked versions.
This came directly from the Era II lesson where in-process mocks passed but
the production database state caused failures.

\subsection{Coverage is a proxy, not a goal}
The push from 17.6\% to 76.6\% coverage in a single day was productive because
the tests were written against real interfaces and uncovered real bugs.
Coverage numbers are meaningless if the tests don't exercise failure paths.

\section{On Velocity}

\subsection{A 54-day gap is a feature, not a failure}
The 54 days between Era II (April 10) and Era III (June 3) with no commits
appear as a gap in the commit history.  The actual work — architectural design,
database schema redesign, multi-role RBAC planning — happened offline.  The
subsequent 52-commit rebuild week was fast precisely because the design was done
before the coding started.

\subsection{Name things for what they are}
The field-name audit sweep (commit \texttt{a0a72bd}) renamed 14 inconsistently-
named fields across the schema and codebase.  Consistent naming prevents class
of bugs that stem from using the wrong field name — bugs that are invisible to
type checkers if the types are compatible.

% ╔══════════════════════════════════════════════════════╗
% ║  CHAPTER 19 — FRONTEND ARCHITECTURE                  ║
% ╚══════════════════════════════════════════════════════╝
\chapter{Frontend Architecture}

\section{Next.js 15 App Router}

The frontend uses Next.js 15 with the App Router — all routes live under
\texttt{app/} with \texttt{page.tsx} components.  Server components are used
where possible for SEO; client components are marked with \texttt{"use client"}
for interactivity.

\section{State Management}

\begin{center}\small
\begin{tabular}{p{3cm}p{4cm}p{6cm}}
\toprule
\rowcolor{headbg}\th{Library} & \th{Scope} & \th{What It Manages} \\
\midrule
Zustand          & Client global     & Auth state, user profile, active role \\
\rowcolor{rowalt}React Query      & Server state      & Restaurant data, orders, menu items, caching \\
React local state & Component        & Form values, UI toggles, modal state \\
\rowcolor{rowalt}Socket.IO client & Real-time events  & Order status, rider location, notifications \\
\bottomrule
\end{tabular}
\end{center}

\section{Page Inventory (44 pages)}

\begin{center}\footnotesize
\begin{longtable}{p{5.5cm}p{7cm}}
\toprule
\rowcolor{headbg}\th{Route} & \th{Purpose} \\
\midrule
\texttt{/}                          & Smart homepage — favourites, trending, search \\
\rowcolor{rowalt}\texttt{/restaurants/[id]}            & Restaurant menu with cart integration \\
\texttt{/checkout}                  & Cashfree drop-in payment \\
\rowcolor{rowalt}\texttt{/checkout/callback}           & Post-payment status (pending state) \\
\texttt{/orders/[id]}               & Order tracking with real-time GPS \\
\rowcolor{rowalt}\texttt{/orders}                      & Order history with re-order \\
\texttt{/favorites}                 & Favourite restaurants \\
\rowcolor{rowalt}\texttt{/profile}                     & Customer profile and addresses \\
\texttt{/auth/login}                & JWT-based login \\
\rowcolor{rowalt}\texttt{/auth/register}               & Multi-step registration \\
\texttt{/onboarding/restaurant}     & Restaurant registration (6,708 cities) \\
\rowcolor{rowalt}\texttt{/onboarding/rider}            & Rider registration \\
\texttt{/shop}                      & Restaurant dashboard: menu, orders \\
\rowcolor{rowalt}\texttt{/shop/menu}                   & Menu CRUD with drag-and-drop \\
\texttt{/shop/orders}               & Order management for restaurant \\
\rowcolor{rowalt}\texttt{/shop/analytics}              & Restaurant revenue analytics \\
\texttt{/delivery/home}             & Rider dashboard: online toggle, current order \\
\rowcolor{rowalt}\texttt{/delivery/earnings}           & Rider earnings and session history \\
\texttt{/admin}                     & Admin landing: quick stats \\
\rowcolor{rowalt}\texttt{/admin/users}                 & User management table \\
\texttt{/admin/restaurants}         & Restaurant management \\
\rowcolor{rowalt}\texttt{/admin/orders}                & Order pipeline view \\
\texttt{/admin/delivery}            & Live rider map (Leaflet) \\
\rowcolor{rowalt}\texttt{/admin/operations}            & Incident list, SLA view \\
\texttt{/admin/analytics}           & Executive analytics (5 sub-views) \\
\rowcolor{rowalt}\texttt{/admin/finance}               & Payments and refunds \\
\texttt{/admin/settings}            & SiteConfig management \\
\rowcolor{rowalt}\texttt{/admin/maintenance}           & Maintenance mode controls \\
\texttt{/admin/coupons}             & Coupon management \\
\rowcolor{rowalt}\texttt{/admin/notifications}         & Broadcast notification management \\
\bottomrule
\end{longtable}
\end{center}

\section{Real-Time Integration}

The Socket.IO client initialises on login and tears down on logout.  A single
connection is shared across all page components via a React context
(\texttt{SocketContext}).  Components subscribe to events via \texttt{useEffect}
hooks that register listeners on mount and remove them on cleanup:

\begin{lstlisting}[language=JavaScript,caption={Socket listener pattern}]
useEffect(() => {
  if (!socket) return;
  const handler = (data) => setOrderStatus(data.status);
  socket.on("order:status-updated", handler);
  return () => socket.off("order:status-updated", handler);
}, [socket]);
\end{lstlisting}

\section{Accessibility (WCAG 2.1 AA)}

Commit \texttt{e4c58a7} addressed WCAG 2.1 AA compliance:

\begin{itemize}[leftmargin=1.4em]
  \item All interactive elements have ARIA labels
  \item Colour contrast meets 4.5:1 ratio for normal text
  \item Focus indicators visible on keyboard navigation
  \item Form inputs have associated \texttt{<label>} elements
  \item Images have descriptive \texttt{alt} text
  \item Screen reader announcements on order status changes
\end{itemize}

% ╔══════════════════════════════════════════════════════╗
% ║  CHAPTER 20 — FUTURE ROADMAP                         ║
% ╚══════════════════════════════════════════════════════╝
\chapter{Future Roadmap}

\section{Near-Term (1–3 months)}

\begin{center}\small
\begin{tabular}{p{4cm}p{3cm}p{5.5cm}}
\toprule
\rowcolor{headbg}\th{Feature} & \th{Priority} & \th{Notes} \\
\midrule
Redis-backed rate limiting & High & Replace in-process; survives redeploys \\
\rowcolor{rowalt}Redis-backed rate limiting & High & Replace in-process; survives redeploys \\
Frontend test coverage  & High    & Target 60\%+ across all 44 pages \\
\rowcolor{rowalt}Image upload (S3/R2)    & Medium  & Replace URL-only model \\
Dead letter queue        & Medium  & Webhook retry with exponential backoff \\
\rowcolor{rowalt}Push notifications       & Medium  & PWA + Firebase for mobile riders \\
\bottomrule
\end{tabular}
\end{center}

\section{Medium-Term (3–6 months)}

\begin{itemize}[leftmargin=1.4em]
  \item \textbf{Multi-tenant}: Each restaurant gets its own subdomain with
        isolated data view; shared infrastructure.
  \item \textbf{Auto-assignment engine}: Replace manual admin assignment with
        proximity-based automatic dispatch using RiderLocation data.
  \item \textbf{Customer mobile app}: React Native with the same Zustand + React
        Query stack; shared service layer.
  \item \textbf{Surge pricing engine}: Dynamic delivery fee based on demand
        (rider count vs open order count) in real time.
  \item \textbf{ML recommendations}: Replace the weighted-composite algorithm
        with a collaborative filtering model trained on the order history dataset.
\end{itemize}

\section{Long-Term (6–12 months)}

\begin{itemize}[leftmargin=1.4em]
  \item \textbf{Operator white-label}: Each restaurant chain deploys a branded
        instance with custom domain, logo, and colour palette.
  \item \textbf{Bank payouts}: Direct payout integration so rider earnings are
        transferred to bank accounts automatically.
  \item \textbf{Restaurant POS integration}: Webhook-based order push to POS
        systems (Petpooja, UrbanPiper) so orders appear in the kitchen display.
  \item \textbf{Subscription tiers}: Restaurant subscription plans with tiered
        commission (0\% for self-hosted, 5\% for platform-managed).
\end{itemize}

% ╔══════════════════════════════════════════════════════╗
% ║  APPENDICES                                          ║
% ╚══════════════════════════════════════════════════════╝
\appendix

\chapter{Complete Commit Timeline}

This appendix lists the five engineering eras with representative commits.
The full 194-commit history is available in the Git log (\texttt{git log --oneline}).

\section{Era I — Genesis (March 26 – April 7)}

\begin{center}\footnotesize
\begin{longtable}{p{2cm}p{10.5cm}}
\toprule
\rowcolor{headbg}\th{Hash} & \th{Message} \\
\midrule
\texttt{cad7cd2} & Phase 1: Production-Ready Auth System with Refresh Tokens \\
\rowcolor{rowalt}\texttt{8aaa7ff} & Phase 2: Cart system with database persistence \\
\texttt{b36822c} & Phase 3: Order System — Cart→Checkout→Order Flow \\
\rowcolor{rowalt}\texttt{b6a92f1} & Cashfree Payment Integration with create, verify, webhook flow \\
\texttt{b8b7fe4} & Phase 2: Real-time order tracking with Socket.IO \\
\rowcolor{rowalt}\texttt{c13d76b} & Phase 3: Admin panel and dashboard \\
\bottomrule
\end{longtable}
\end{center}

\section{Era II — Production Crisis (April 7–10)}

\begin{center}\footnotesize
\begin{longtable}{p{2cm}p{10.5cm}}
\toprule
\rowcolor{headbg}\th{Hash} & \th{Message} \\
\midrule
\texttt{---} & ISSUE 1–3: Production deployment fixes \\
\rowcolor{rowalt}\texttt{---} & ISSUE 5: Add transaction locking for race condition prevention \\
\texttt{---} & ISSUE 6: Fix order status override with conditional update \\
\rowcolor{rowalt}\texttt{---} & ISSUE 7: Optimize cron job with single atomic updateMany \\
\texttt{---} & ISSUE 8: Add database indexes for Order model queries \\
\rowcolor{rowalt}\texttt{---} & FINAL PRODUCTION GAPS: Fix 3 critical vulnerabilities \\
\bottomrule
\end{longtable}
\end{center}

\section{Era III — The Great Rebuild (June 3–7)}

\begin{center}\footnotesize
\begin{longtable}{p{2cm}p{10.5cm}}
\toprule
\rowcolor{headbg}\th{Hash} & \th{Message} \\
\midrule
\texttt{bef7fcf} & Wait for Zustand hydration before auth redirect \\
\rowcolor{rowalt}\texttt{15d1a98} & Don't clear auth state on cross-origin 401 \\
\texttt{24f47a6} & Use Authorization header to bypass third-party cookie blocking \\
\rowcolor{rowalt}\texttt{17a035e} & Allow Authorization header in preflight CORS response \\
\texttt{02c1bf8} & Shop menu page with drag-and-drop reordering \\
\rowcolor{rowalt}\texttt{4216492} & Add 6,708 Indian cities dataset for restaurant onboarding \\
\bottomrule
\end{longtable}
\end{center}

\section{Era IV — Product Maturity (June 10–13)}

\begin{center}\footnotesize
\begin{longtable}{p{2cm}p{10.5cm}}
\toprule
\rowcolor{headbg}\th{Hash} & \th{Message} \\
\midrule
\texttt{bc9090d} & Bootstrap admin endpoint (zero-admin guard) \\
\rowcolor{rowalt}\texttt{59ff07a} & Recommendations engine with personalised weighting \\
\texttt{eecf1c9} & IDOR closure — ownership gates on all resource endpoints \\
\rowcolor{rowalt}\texttt{e4c58a7} & WCAG 2.1 AA accessibility fixes \\
\texttt{ed0b03d} & Backend test coverage push — begin 927-test sprint \\
\rowcolor{rowalt}\texttt{cd7462b} & Backend test coverage sprint complete — 76.6\% \\
\texttt{f065ac6} & GitHub Actions CI/CD pipeline \\
\rowcolor{rowalt}\texttt{a0a72bd} & Field-name audit sweep — 14 inconsistencies resolved \\
\bottomrule
\end{longtable}
\end{center}

\section{Era V — Operations Excellence (June 14)}

\begin{center}\footnotesize
\begin{longtable}{p{2cm}p{10.5cm}}
\toprule
\rowcolor{headbg}\th{Hash} & \th{Message} \\
\midrule
\texttt{0e36a14} & feat(ops): Ops Intelligence Wave A — incident engine \\
\rowcolor{rowalt}\texttt{d33f5cb} & feat(ops): Ops Intelligence Wave B — escalation + scoring \\
\texttt{51f0a80} & feat(ops): Ops Intelligence Wave C — executive analytics \\
\rowcolor{rowalt}\texttt{02c7962} & feat(security): OWASP hardening sprint part 1 \\
\texttt{0f21179} & feat(security): Security sprint part 2 — audit trail + regression suite \\
\rowcolor{rowalt}\texttt{7fe0497} & feat(scale): Distributed locks for cron jobs + Redis-backed SiteConfig cache \\
\texttt{2d606b3} & feat(obs): Production observability hardening — Phases 1--3 \\
\rowcolor{rowalt}\texttt{987e9b5} & feat(obs): Phase 4 — route cron job errors to Sentry \\
\texttt{37593ee} & fix(delivery): extract riderTracker service, eliminate duplicate GPS watcher \\
\rowcolor{rowalt}\texttt{00bfc3d} & fix(payment): correct Cashfree refund SDK method name \\
\texttt{1395454} & docs(obs): GhostKitchen\_Production\_Observability\_Report.tex \\
\rowcolor{rowalt}\texttt{f8516e5} & docs(audit): Master Project Lifecycle Audit \& Certification Report \\
\bottomrule
\end{longtable}
\end{center}

\chapter{Prisma Schema Reference}

The complete schema is defined in \texttt{food-delivery-backend/prisma/schema.prisma}.
This appendix documents the 20 models and their field counts.

\begin{center}\small
\begin{longtable}{p{3.5cm}rp{8cm}}
\toprule
\rowcolor{headbg}\th{Model} & \th{Fields} & \th{Key Relationships} \\
\midrule
User          & 31 & RefreshToken[], Order[], Restaurant, RiderLocation, AuditLog[], Address[], Favorite[], RiderSession[] \\
\rowcolor{rowalt}Restaurant    & 23 & User(owner), MenuItem[], Order[], Review[], Coupon[], Favorite[] \\
Order         & 24 & User(customer), Restaurant, User(agent), CartItem[], Payment, Refund \\
\rowcolor{rowalt}Payment       & 12 & Order, PaymentWebhook[], Refund[] \\
MenuItem      & 15 & Restaurant, CartItem[] \\
\rowcolor{rowalt}CartItem      & 8  & User, MenuItem, Order \\
Coupon        & 12 & Restaurant, Order[] \\
\rowcolor{rowalt}Review        & 9  & User, Restaurant, Order \\
Address       & 14 & User \\
\rowcolor{rowalt}Favorite      & 6  & User, Restaurant (unique constraint) \\
Notification  & 9  & User \\
\rowcolor{rowalt}AuditLog      & 9  & User \\
RefreshToken  & 7  & User \\
\rowcolor{rowalt}RiderLocation & 8  & User (unique on riderId) \\
RiderSession  & 9  & User \\
\rowcolor{rowalt}Incident      & 20 & IncidentEvent[] \\
IncidentEvent & 9  & Incident \\
\rowcolor{rowalt}PaymentWebhook & 7  & Payment \\
Refund        & 9  & Order, Payment \\
\rowcolor{rowalt}SiteConfig    & 17 & Singleton (no FK) \\
\bottomrule
\end{longtable}
\end{center}

\chapter{API Endpoint Reference}

\begin{center}\footnotesize
\begin{longtable}{p{1.5cm}p{5cm}p{6cm}}
\toprule
\rowcolor{headbg}\th{Method} & \th{Path} & \th{Description} \\
\midrule
\multicolumn{3}{l}{\textbf{Auth — /api/auth}} \\
\midrule
POST & /register               & Create account (email, password, name) \\
\rowcolor{rowalt}POST & /login                  & Email + password → access + refresh tokens \\
GET  & /me                     & Validate token; auto-heal restaurantId if missing \\
\rowcolor{rowalt}POST & /refresh                & Rotate refresh token; return new access token \\
POST & /logout                 & Revoke current refresh token \\
\rowcolor{rowalt}POST & /logout-all             & Revoke all refresh tokens for user \\
\midrule
\multicolumn{3}{l}{\textbf{Role — /api/role}} \\
\midrule
POST & /switch                 & Switch active role; re-issue token \\
\rowcolor{rowalt}POST & /register-restaurant    & Add RESTAURANT role + create Restaurant row \\
POST & /register-rider         & Add DELIVERY role + create RiderLocation row \\
\rowcolor{rowalt}POST & /add-restaurant         & Admin: add restaurant to existing user \\
GET  & /my-restaurant          & Get calling user's restaurant \\
\midrule
\multicolumn{3}{l}{\textbf{Restaurant — /api/restaurants}} \\
\midrule
GET  & /                       & List restaurants (filter, paginate, search) \\
\rowcolor{rowalt}GET  & /:id                    & Get restaurant with menu \\
POST & /                       & Create restaurant (RESTAURANT role) \\
\rowcolor{rowalt}PUT  & /:id                    & Update restaurant (owner only) \\
PATCH & /:id/set-status        & Toggle open/closed + statusNote \\
\rowcolor{rowalt}GET  & /:id/analytics          & Revenue, orders, top items \\
POST & /:id/menu               & Add menu item \\
\rowcolor{rowalt}PUT  & /:id/menu/:itemId       & Update menu item \\
DELETE & /:id/menu/:itemId     & Delete menu item \\
\rowcolor{rowalt}POST & /trending               & Top 10 by recent order volume \\
GET  & /recommendations        & Personalised restaurant list \\
\midrule
\multicolumn{3}{l}{\textbf{Orders — /api/orders}} \\
\midrule
GET  & /                       & List orders for calling user's role \\
\rowcolor{rowalt}GET  & /:id                    & Get order (ownership check) \\
POST & /                       & Place order (transactional) \\
\rowcolor{rowalt}PATCH & /:id/status            & Transition order status (role-gated) \\
POST & /:id/reorder            & Clone order to new cart \\
\midrule
\multicolumn{3}{l}{\textbf{Payments — /api/payments}} \\
\midrule
POST & /create-order           & Cashfree PGCreateOrder; return session\_id \\
\rowcolor{rowalt}GET  & /verify/:orderId        & Check payment status from Cashfree \\
POST & /webhook                & HMAC-verified Cashfree webhook handler \\
\rowcolor{rowalt}GET  & /history                & Payment history for calling user \\
POST & /my-refunds             & Customer-initiated refund request \\
\rowcolor{rowalt}GET  & /my-refunds             & List customer's refunds \\
POST & /refunds                & Admin: initiate refund \\
\rowcolor{rowalt}GET  & /refunds                & Admin: list all refunds \\
PATCH & /refunds/:id/sync      & Sync refund status from Cashfree \\
\midrule
\multicolumn{3}{l}{\textbf{Delivery — /api/delivery}} \\
\midrule
PATCH & /status                & Toggle online/offline + GPS coordinates \\
\rowcolor{rowalt}POST & /location              & Update GPS position (every 5s) \\
GET  & /earnings               & Earnings + session analytics \\
\rowcolor{rowalt}GET  & /current-order         & Active delivery order \\
\midrule
\multicolumn{3}{l}{\textbf{Admin — /api/admin (selection)}} \\
\midrule
POST & /bootstrap-admin        & Create first admin (zero-admin guard) \\
\rowcolor{rowalt}GET  & /users                  & List all users (paginated) \\
PATCH & /users/:id/suspend     & Suspend user + revoke tokens \\
\rowcolor{rowalt}PATCH & /users/:id/block       & Block user + revoke tokens \\
GET  & /orders                 & All orders with filters \\
\rowcolor{rowalt}POST & /orders/:id/assign-delivery & Assign rider to order \\
GET  & /analytics/revenue      & Platform-wide revenue analytics \\
\rowcolor{rowalt}GET  & /audit-log              & Immutable admin action log \\
GET  & /audit-log/export       & CSV export of audit log \\
\midrule
\multicolumn{3}{l}{\textbf{Health — /health}} \\
\midrule
GET  & /health                 & Readiness: DB + Redis status \\
\rowcolor{rowalt}GET  & /health/extended        & Full diagnostic with cron, socket, maintenance \\
\bottomrule
\end{longtable}
\end{center}

\chapter{Security Quick Reference}

\section{RBAC Matrix}

\begin{center}\small
\begin{tabular}{p{4cm}cccc}
\toprule
\rowcolor{headbg}\th{Operation} & \th{CUSTOMER} & \th{RESTAURANT} & \th{DELIVERY} & \th{ADMIN} \\
\midrule
Place order          & \yes{} & \no{}  & \no{}  & \yes{} \\
\rowcolor{rowalt}Confirm/prepare order & \no{}  & \yes{} & \no{}  & \yes{} \\
Accept delivery      & \no{}  & \no{}  & \yes{} & \yes{} \\
\rowcolor{rowalt}Manage menu          & \no{}  & \yes{} & \no{}  & \yes{} \\
View own orders      & \yes{} & \yes{} & \yes{} & \yes{} \\
\rowcolor{rowalt}View all orders      & \no{}  & \no{}  & \no{}  & \yes{} \\
Manage users         & \no{}  & \no{}  & \no{}  & \yes{} \\
\rowcolor{rowalt}Access exec analytics & \no{}  & \no{}  & \no{}  & \yes{} \\
Initiate refund      & \yes{} & \no{}  & \no{}  & \yes{} \\
\rowcolor{rowalt}Manage coupons       & \no{}  & \yes{} & \no{}  & \yes{} \\
\bottomrule
\end{tabular}
\end{center}

\section{Rate Limits}

\begin{center}\small
\begin{tabular}{p{4cm}p{2.5cm}p{6cm}}
\toprule
\rowcolor{headbg}\th{Endpoint Group} & \th{Limit} & \th{Rationale} \\
\midrule
General (all)    & 200 / 15 min & Baseline traffic protection \\
\rowcolor{rowalt}Auth endpoints   & 10 / 15 min  & Brute-force protection on login/register \\
Payment creation & 20 / 1 min   & Fraud prevention \\
\rowcolor{rowalt}Admin operations & 100 / 15 min & Admin workflows need higher throughput \\
\bottomrule
\end{tabular}
\end{center}

\chapter{Glossary}

\begin{description}[leftmargin=4.5cm,style=nextline]
  \item[CUID2] Collision-resistant unique identifier v2 — used as primary keys.
  \item[Distributed Lock] Redis SETNX pattern ensuring only one instance executes a critical section.
  \item[HMAC-SHA256] Hash-based Message Authentication Code — used to verify Cashfree webhooks.
  \item[IDOR] Insecure Direct Object Reference — accessing resources via guessable IDs without ownership check.
  \item[Idempotent DDL] Database schema changes that are safe to apply multiple times (IF NOT EXISTS).
  \item[Paise] Smallest Indian currency unit: ₹1 = 100 paise. All monetary values stored in paise.
  \item[Prisma] TypeScript ORM for Node.js providing type-safe database access and migration management.
  \item[Refresh Token Rotation] Invalidating the old refresh token on every use and issuing a new one.
  \item[Socket Room] A Socket.IO channel scoped to a specific resource (order, restaurant, or rider).
  \item[SiteConfig] Singleton database table holding platform-wide settings (maintenance mode, fees, flags).
  \item[Three-Layer Cache] In-process (5s) → Redis (30s) → PostgreSQL. Used for SiteConfig.
  \item[Upstash] Redis-compatible serverless data store accessible via HTTP REST.
  \item[WCAG 2.1 AA] Web Content Accessibility Guidelines Level AA — the accessibility standard implemented.
  \item[Zod] TypeScript-first schema validation library used for all API request validation.
\end{description}

\backmatter

\chapter*{Colophon}
\addcontentsline{toc}{chapter}{Colophon}

\begin{center}
\begin{tabular}{p{4cm}p{9cm}}
\toprule
\textbf{Document}      & GhostKitchen Complete Engineering Journey \\
\textbf{Version}       & 1.0 \\
\textbf{Date}          & June 2026 \\
\textbf{Format}        & \LaTeX\ (book class, 11pt, A4) \\
\textbf{Typeset with}  & pdflatex + TikZ + pgfplots + tcolorbox \\
\textbf{Data sources}  & git log, prisma/schema.prisma, package.json, CI coverage reports \\
\textbf{LOC verified}  & June 14, 2026 (cloc) \\
\textbf{Commit count}  & 194 (March 26 – June 14, 2026) \\
\bottomrule
\end{tabular}

\vspace{2cm}

\begin{tcolorbox}[colback=brand!6,colframe=brand,width=12cm]
\centering
\textbf{GhostKitchen Engineering}\\[4pt]
\small Production-ready platform built in 81 days.\\
22,767 + 27,188 source lines. 1,036 tests. 194 commits.
\end{tcolorbox}
\end{center}

\end{document}

